\chapter{Sectional Curvature Comparison I}
\label{chap:pet-sectional-curvature-comparison-i}

Faithful transcription of Petersen, \emph{Riemannian Geometry} (GTM 171, 3rd ed.),
Chapter 6. Having classified the complete spaces of constant curvature, the chapter
compares manifolds of variable curvature to these model spaces. The first global
result is the Hadamard--Cartan theorem: a simply connected complete manifold with
$\sec\le0$ is diffeomorphic to $\mathbb{R}^n$. In positive curvature there are strong
topological restrictions — the Bonnet--Myers diameter bound and Synge's theorem
that an orientable even-dimensional manifold with $\sec>0$ is simply connected — and
the chapter culminates in the classical quarter-pinched sphere theorem of Rauch,
Berger, and Klingenberg. The technical engine throughout is the differentiation of
vector fields along a single curve, which is used to define parallel fields and
Jacobi fields and to establish Synge's second variation formula of the energy
functional; basic Riccati and Rauch comparison estimates for Hessians of distance
functions are then used to bound injectivity, conjugate, and convexity radii.

\section{The Connection Along Curves}

\begin{lemma}[the chart--connection bridge]
  \label{lem:chart-connection-bridge}
  \lean{PetersenLib.leviCivita_chartFrame_christoffel,
    PetersenLib.exists_chartFrame_leviCivita_christoffel_nhds}
  \leanok
  Infrastructure, not a numbered result of the source; recorded because the
  formalization needs it and Petersen takes it for granted.

  Around any $p$ in a chart around $\alpha$ there is a local frame $Z_1,\dots,Z_n$
  and a neighbourhood $U\ni p$ on which the Levi-Civita connection of $g$ is
  computed by the Christoffel symbols of the second kind:
  \[
    \nabla_{Z_i}Z_j=\sum_m\Gamma^m_{ij}\,Z_m,\qquad
    \Gamma^m_{ij}=\tfrac12\sum_lG^{ml}\big(\partial_iG_{lj}+\partial_jG_{li}
      -\partial_lG_{ij}\big).
  \]
  Here $\nabla$ is the abstract (Koszul) connection through which the curvature
  tensor $R$ and $\sec$ are defined, while $\Gamma$ is the chart symbol through
  which geodesics, $\exp$, and the variational calculus are defined. Petersen
  moves between the two without comment; in the formalization they are distinct
  objects, and every result of this chapter that couples curvature to geodesics
  must identify them.

  \emph{Remark on how this is actually used.} The identification that
  \cref{lem:pet-ch6-third-partial-commutator} consumes is the sharper
  \cref{lem:curvature-chart-bridge}, which reaches all the way to $R$ rather than
  stopping at $\nabla$, and which is obtained from
  \cref{prop:pet-ch3-curvature-coordinate-formula} rather than from this lemma.
  This lemma remains the general-purpose statement of the identification (and is
  strictly more general than its Chapter 3 counterpart, being stated for an
  arbitrary chart centre $\alpha$ rather than only for $\alpha=p$).

  The proof also needs germ-locality of $\nabla$ (that $\nabla_vX$ depends only on
  the germ of $X$), which the source likewise leaves implicit.
\end{lemma}

\begin{lemma}[the curvature--chart bridge]
  \label{lem:curvature-chart-bridge}
  \lean{PetersenLib.chartCurvatureContraction2_eq_neg_curvatureTensorAt}
  \leanok
  \uses{prop:pet-ch3-curvature-coordinate-formula}
  Infrastructure, not a numbered result of the source; recorded because the
  formalization needs it and Petersen takes it for granted.

  Let $R_{\mathrm{chart}}(X,Y)Z=\sum_l\big(\sum_{i,j,k}R^l{}_{ijk}(y)X^iY^jZ^k\big)e_l$
  be the curvature contraction written directly in the Christoffel symbols of the
  chart at $p$, with
  \[
    R^l{}_{ijk}=\partial_j\Gamma^l_{ik}-\partial_i\Gamma^l_{jk}
      +\sum_s\big(\Gamma^s_{ik}\Gamma^l_{js}-\Gamma^s_{jk}\Gamma^l_{is}\big).
  \]
  Then, at the centre of that chart,
  \[
    R_{\mathrm{chart}}(X,Y)Z=-R(X,Y)Z,
  \]
  with $R$ the abstract (Koszul) curvature tensor of \cref{def:pet-ch3-curvature-tensor}.

  Two remarks on the shape of the statement.

  \emph{The sign.} $R_{\mathrm{chart}}$ is written in do Carmo's convention
  $R(X,Y)=\nabla_Y\nabla_X-\nabla_X\nabla_Y+\nabla_{[X,Y]}$, the negative of the
  convention of \cref{def:pet-ch3-curvature-tensor}; the two coefficient
  expressions are exact negations with the same index slots. In
  \cref{lem:pet-ch6-third-partial-commutator} this sign cancels against the
  opposite parameter order in which the coordinate commutator is proved, so
  Petersen's identity holds with Petersen's $R$ and no stray sign.

  \emph{Why the chart centre suffices.}
  \cref{prop:pet-ch3-curvature-coordinate-formula} is stated on the diagonal —
  chart centre $=$ evaluation point. This is no restriction for
  \cref{lem:pet-ch6-third-partial-commutator}, whose conclusion is pointwise at
  $p=c(u_0,s_0,t_0)$: read the map in the chart at that very point. The step that
  makes the diagonal case collapse is that at the centre of its own chart the
  chart frame $\partial_i|_p$ \emph{is} the model basis $e_i$, so that the
  chart-frame expansion of a tangent vector at $p$ is the ordinary basis expansion
  of $E$ and the coordinate-to-abstract comparison map is the identity.
\end{lemma}

\begin{definition}[covariant derivative along a curve]
  \label{def:pet-ch6-vector-field-along-curve}
  \lean{PetersenLib.derivAlongCurve,PetersenLib.covariantDerivCoord_transfer}
  \leanok
  For a curve $c:I\to M$, a \emph{vector field along $c$} is a map $V:I\to TM$
  with $V(t)\in T_{c(t)}M$ for all $t\in I$. Thinking of $V$ as the variational
  field of a variation $\bar c:(-\varepsilon,\varepsilon)\times I\to M$ with
  $\bar c(0,t)=c(t)$, $\frac{\partial\bar c}{\partial s}(0,t)=V(t)$, the
  \emph{covariant derivative of $V$ along $c$} is defined by
  $\dot V(t)=\frac{d}{dt}V(t)=\frac{\partial^2\bar c}{\partial t\,\partial s}(0,t)$.
  In local coordinates, writing $V(t)=V^k(t)\partial_k$ and
  $\bar c^k(s,t)=\tfrac{\partial\bar c^k}{\partial s}(0,t)$, one computes
  \[
    \frac{\partial^2\bar c}{\partial t\,\partial s}(0,t)
    =\dot V^k(t)\partial_k+V^i(t)\dot c^j(t)\Gamma^k_{ij}\partial_k,
  \]
  which is manifestly independent of the choice of variation and of the
  coordinate system, and agrees with $\nabla_{\dot c(t_0)}X$ whenever
  $V(t)=X_{c(t)}$ for a vector field $X$ near $c(t_0)$. In general $\dot V(t_0)$
  need not vanish when $\dot c(t_0)=0$ (e.g.\ for a constant curve $c$, where $V$
  is simply a curve in the fixed space $T_{c(t_0)}M$). Selecting a two-parameter
  variation as needed for a second field $W$ along $c$, the local formula also
  yields the product rule $\frac{d}{dt}g(V,W)=g(\dot V,W)+g(V,\dot W)$ and
  linearity $\frac{d}{dt}(V+W)=\dot V+\dot W$,
  $\frac{d}{dt}(\lambda V)=\dot\lambda V+\lambda\dot V$ for $\lambda:I\to\mathbb{R}$.
  When $M\subset\tilde M$, the derivative along $c$ in $M$ is obtained by
  computing $\dot V\in T\tilde M$ and orthogonally projecting onto $TM$.
  \source{petersen:page-0248,petersen:page-0249}
\end{definition}

\begin{definition}[third and higher partial derivatives]
  \label{def:pet-ch6-higher-order-partials}
  \lean{PetersenLib.surfaceCovariantDeriv,PetersenLib.thirdPartialDerivative}
  \leanok
  \uses{def:pet-ch6-vector-field-along-curve}
  Using \cref{def:pet-ch6-vector-field-along-curve}, higher mixed partials of a
  multi-parameter map $c$ are defined recursively: to compute
  $\frac{\partial^3c}{\partial s\,\partial t\,\partial u}(s_0,t_0,u_0)$, set
  $V(s)=\frac{\partial^2c}{\partial t\,\partial u}(s,t_0,u_0)$ and define the
  third partial as $\frac{dV}{ds}(s_0)$. Because this recursive definition fixes
  a specific order of differentiation, there is no a priori reason for third
  partials to be symmetric in all three variables; it remains true, however,
  that $\frac{\partial^3c}{\partial s\,\partial t\,\partial u}
  =\frac{\partial^3c}{\partial s\,\partial u\,\partial t}$ (the last two
  variables, which enter only through the already-defined second mixed
  partial, may be swapped).
  \source{petersen:page-0249}
\end{definition}

\begin{example}[Example 6.1.1 — mixed partials on $S^2(1)$]
  \label{ex:pet-ch6-sphere-third-partials}
  \lean{PetersenLib.sphereThirdPartialsExample}
  \uses{def:pet-ch6-higher-order-partials}
  Parametrize $S^2(1)\subset\mathbb{R}^3$ as a surface of revolution about the
  $x$-axis by $c(t,\theta)=(\cos t,\sin t\cos\theta,\sin t\sin\theta)$, and
  compute all derivatives in $\mathbb{R}^3$ before projecting onto $S^2(1)$. The
  curves $t\mapsto c(t,\theta)$ are geodesics, since
  $\partial_tc=(-\sin t,\cos t\cos\theta,\cos t\sin\theta)\in TS^2(1)$ while
  $\partial_t^2c=(-\cos t,-\sin t\cos\theta,-\sin t\sin\theta)=-c$ is normal to
  $S^2(1)$ and so projects to zero. Next,
  $\partial_\theta\partial_tc=(0,-\cos t\sin\theta,\cos t\cos\theta)$ is already
  tangent to $S^2(1)$, hence equals the intrinsic mixed partial. Finally
  \[
    \partial_t\partial_\theta\partial_tc=(0,\sin t\sin\theta,-\sin t\cos\theta)
    =\partial_\theta\partial_t^2c
  \]
  in $\mathbb{R}^3$, and both are tangent to $S^2(1)$; the first is
  $\partial_t\partial_\theta\partial_tc$ as computed intrinsically in $S^2(1)$,
  while the second has no intrinsic meaning as written, since $\partial_t^2c$
  must first be projected to $S^2(1)$ (where it vanishes) before differentiating
  in $\theta$. Hence intrinsically
  $\partial_\theta\partial_t^2c=0\ne\partial_t\partial_\theta\partial_tc$ along
  $t=0$, illustrating that third partials computed with different orderings of
  differentiation need not agree. Along the equator $t=0$,
  $\partial_\theta\partial_tc=0$.
  \source{petersen:page-0250}
\end{example}

\begin{lemma}[Lemma 6.1.2 — third partials and curvature]
  \label{lem:pet-ch6-third-partial-commutator}
  \lean{PetersenLib.thirdPartialCurvatureFormula}
  \leanok
  \uses{def:pet-ch6-higher-order-partials, lem:curvature-chart-bridge}
  For a map $c(u,s,t)$ with values in $M$,
  \[
    \frac{\partial^3c}{\partial u\,\partial s\,\partial t}
    -\frac{\partial^3c}{\partial s\,\partial u\,\partial t}
    =R\!\left(\frac{\partial c}{\partial u},\frac{\partial c}{\partial s}\right)
    \frac{\partial c}{\partial t}.
  \]

  \emph{Formalization note.} The Lean statement is the identity for the
  $(u,s)$-surface with an \emph{arbitrary} $C^2$ field $V$ along it, of which
  Petersen's statement is the case $V=\frac{\partial c}{\partial t}$: the $t$-slot
  enters only through $V$, and the general form is what the Jacobi equation
  (\cref{def:pet-ch6-jacobi-field}) actually consumes. The surface is read in the
  chart centred at the base point $p=c(u_0,s_0,t_0)$, which by
  \cref{lem:curvature-chart-bridge} costs no generality, and the regularity Petersen
  leaves implicit is explicit.
  \source{petersen:page-0251}
\end{lemma}
\begin{proof}
  Fix $p=c(u,s,t)$ and choose coordinates with $g_{ij}|_p=\delta_{ij}$,
  $\Gamma^k_{ij}|_p=0$, so that $\frac{\partial}{\partial u}(\partial_i)|_p=0$.
  Then at $p$,
  \[
    \frac{\partial^3c}{\partial u\,\partial s\,\partial t}
    =\frac{\partial}{\partial u}\!\left(\frac{\partial^2c^l}{\partial s\,\partial t}\partial_l
    +\frac{\partial c^i}{\partial s}\frac{\partial c^j}{\partial t}\Gamma^k_{ij}\partial_k\right)
    =\frac{\partial^3c^l}{\partial u\,\partial s\,\partial t}\partial_l
    +\frac{\partial c^i}{\partial s}\frac{\partial c^j}{\partial u}
    \frac{\partial}{\partial t}(\Gamma^k_{ij})\partial_k,
  \]
  using $\frac{\partial\Gamma^k_{ij}}{\partial u}
  =\frac{\partial c^l}{\partial u}\partial_l\Gamma^k_{ij}$ and swapping the roles of
  $j,u$ in the chain rule term against $\frac{\partial}{\partial u}=\frac{\partial
  c^j}{\partial u}\partial_j$ acting on $\Gamma^k_{ij}$ (this uses only that
  $\Gamma^k_{ij}|_p=0$, so the second-derivative-of-$\Gamma$ terms drop and only
  the first derivative of $\Gamma$ along $u$ survives). Similarly
  \[
    \frac{\partial^3c}{\partial s\,\partial u\,\partial t}
    =\frac{\partial^3c^l}{\partial s\,\partial u\,\partial t}\partial_l
    +\frac{\partial c^i}{\partial s}\frac{\partial c^j}{\partial u}
    \frac{\partial c^k}{\partial t}(\partial_k\Gamma^l_{ij})\partial_l.
  \]
  Since ordinary third partials of the coordinate functions $c^l$ are symmetric
  in $u,s$, subtracting and relabelling the dummy index in the first expression
  gives
  \[
    \frac{\partial^3c}{\partial u\,\partial s\,\partial t}\bigg|_p
    -\frac{\partial^3c}{\partial s\,\partial u\,\partial t}\bigg|_p
    =\frac{\partial c^i}{\partial s}\frac{\partial c^j}{\partial u}
    \frac{\partial c^k}{\partial t}(\partial_k\Gamma^l_{ij}-\partial_j\Gamma^l_{ik})\partial_l
    =\frac{\partial c^i}{\partial s}\frac{\partial c^j}{\partial u}
    \frac{\partial c^k}{\partial t}R^l_{kji}\partial_l,
  \]
  using the coordinate formula for the curvature tensor from section 3.1.6 of the
  source (at a point where $\Gamma^k_{ij}|_p=0$, $R^l_{kji}=\partial_k\Gamma^l_{ij}
  -\partial_j\Gamma^l_{ik}$). This is exactly
  $R\!\left(\frac{\partial c}{\partial u},\frac{\partial c}{\partial s}\right)
  \frac{\partial c}{\partial t}$.

  \emph{Formalization note.} The formal proof is this same coordinate computation,
  but carried out without the normalization $\Gamma^k_{ij}|_p=0$: rather than
  choosing coordinates that kill the quadratic terms, it keeps them and lets them
  cancel, which is the content of do Carmo's Chapter 4 Lemma 4.1 — the identity
  above with $(s,t)$ for $(u,s)$ and $V=\frac{\partial c}{\partial t}$. That
  computation is reused verbatim; the only work is
  \cref{lem:curvature-chart-bridge}, which converts its Christoffel-level
  right-hand side into $R$. Petersen's normalization is thus a convenience of
  exposition rather than a step the formalization needs.
  \source{petersen:page-0251}
  \leanok
\end{proof}

\begin{definition}[parallel field along a curve]
  \label{def:pet-ch6-parallel-field}
  \lean{PetersenLib.IsParallelAlong}
  \leanok
  \uses{def:pet-ch6-vector-field-along-curve}
  A vector field $V$ along $c$ is \emph{parallel along $c$} if $\dot V\equiv0$.
  The tangent field $\dot c$ along a geodesic is parallel, and (by
  \cref{ex:pet-ch6-sphere-third-partials}) the unit field perpendicular to a
  great circle in $S^2(1)$ is parallel. If $V,W$ are parallel along $c$ then
  $g(V,W)$ is constant along $c$ by the product rule, so parallel fields neither
  change length nor relative angle, generalizing parallel translation in
  Euclidean space.
  \source{petersen:page-0252}
\end{definition}

\begin{theorem}[Theorem 6.1.3 — existence and uniqueness of parallel fields]
  \label{thm:pet-ch6-parallel-field-existence-uniqueness}
  \lean{PetersenLib.parallelField_existence_uniqueness_global,
    PetersenLib.exists_isParallelSolOn_Ioo_of_openHyp}
  \leanok
  \uses{def:pet-ch6-parallel-field,thm:pet-ch6-parallel-field-chart-local}
  If $t_0\in I$ and $v\in T_{c(t_0)}M$, there is a unique parallel field $V(t)$
  defined on all of $I$ with $V(t_0)=v$.

  Global existence on the whole of $I$ — not merely locally — is the substance of
  the theorem: it is what distinguishes the \emph{linear} parallel-transport
  equation from the nonlinear geodesic equation, whose solutions exist only
  locally. The formalization assumes only that $c$ is $C^2$ on $I$; no
  chart-containment hypothesis is imposed, and
  \texttt{exists\_isParallelSolOn\_Ioo\_of\_openHyp} takes Petersen's hypothesis
  verbatim, with $c$ regular only on the \emph{open} interval.
  \source{petersen:page-0252}
\end{theorem}
\begin{proof}
  \uses{thm:pet-ch6-parallel-field-chart-local,def:pet-ch6-vector-field-along-curve}
  \leanok
  Petersen solves the linear system $\dot V^j=-\sum_i\alpha^j_iV^i$ in a frame.
  We argue by gluing the chart-local solutions of
  \cref{thm:pet-ch6-parallel-field-chart-local}. Fix a compact
  $[a,b]\subseteq I$ containing $t_0$. By compactness of $c([a,b])$ and a
  Lebesgue-number argument for a cover of it by chart sources, there is a
  partition $a=u_0<\dots<u_n=b$ with each $c([u_i,u_{i+1}])$ inside a single
  chart source; on each piece the Christoffel contraction is continuous and
  bounded, so \cref{thm:pet-ch6-parallel-field-chart-local} applies. Walking
  right from $t_0$ and then left, the pieces are glued at their shared
  endpoints: the local uniqueness forces agreement there, and chart-change
  covariance (\texttt{covariantDerivCoord\_transfer},
  \cref{def:pet-ch6-vector-field-along-curve}) makes ``parallel'' independent of
  the chart, so the glued field is parallel \emph{across} each junction and not
  merely on the interiors. Uniqueness on $[a,b]$ follows by inducting the local
  uniqueness along the partition. Finally, an arbitrary open $I=(a,b)$ is
  exhausted by compacta $[A_n,B_n]\uparrow(a,b)$ all containing $t_0$; uniqueness
  on the smaller window makes the solutions cohere, and the coherent limit is the
  field on all of $I$.
  \source{petersen:page-0252,petersen:page-0253}
\end{proof}

\begin{theorem}[Theorem 6.1.3, chart-local form]
  \label{thm:pet-ch6-parallel-field-chart-local}
  \lean{PetersenLib.parallelField_existence_uniqueness}
  \leanok
  \uses{def:pet-ch6-parallel-field}
  Let $c$ be a curve whose restriction to $(a,b)$ is differentiable in a chart
  around $\alpha$ and stays in that chart's source, with the associated
  Christoffel contraction continuous and bounded on $[a,b]$. Then for $t_0\in(a,b)$
  and $v\in T_{c(t_0)}M$ there is a vector field $V$ along $c$, parallel on
  $(a,b)$, with $V(t_0)=v$, and any parallel field $W$ along $c|_{(a,b)}$ with
  $W(t_0)=v$ agrees with $V$ on $(a,b)$.
  \source{petersen:page-0252}
\end{theorem}
\begin{proof}
  \uses{thm:pet-ch6-parallel-field-existence-uniqueness}
  Choose vector fields $E_1(t),\dots,E_n(t)$ along $c$ forming a basis for
  $T_{c(t)}M$ for every $t\in I$, and write $\nabla_{\dot c}E_i=\sum_j
  \alpha_i^j(t)E_j(t)$. For $V(t)=V^i(t)E_i(t)$,
  \[
    \dot V=\sum_i\dot V^i(t)E_i(t)+\sum_iV^i(t)\nabla_{\dot c}E_i
    =\sum_j\Big(\dot V^j(t)+\sum_iV^i(t)\alpha_i^j(t)\Big)E_j(t),
  \]
  so $V$ is parallel iff $V^1,\dots,V^n$ solve the linear first-order system
  $\dot V^j(t)=-\sum_i\alpha_i^j(t)V^i(t)$, $j=1,\dots,n$. Linear systems of
  ordinary differential equations have, for any prescribed initial values
  $V^1(t_0),\dots,V^n(t_0)$, a unique solution defined on the whole interval
  $I$ (this global existence, in contrast to the nonlinear geodesic equation,
  is a standard consequence of linearity of the equations).
  \source{petersen:page-0252,petersen:page-0253}
\end{proof}

\begin{definition}[Jacobi field]
  \label{def:pet-ch6-jacobi-field}
  \lean{PetersenLib.IsJacobiFieldAlong,PetersenLib.jacobiEquation}
  \leanok
  \uses{def:pet-ch6-vector-field-along-curve, lem:pet-ch6-third-partial-commutator}
  Let $\bar c(s,t)$ be a \emph{geodesic variation} of a geodesic $c(t)=\bar
  c(0,t)$, i.e.\ $t\mapsto\bar c(s,t)$ is a geodesic for every $s$. By
  \cref{lem:pet-ch6-third-partial-commutator} applied with $u=t$ (using that
  $\frac{\partial^2\bar c}{\partial t^2}\equiv0$),
  \[
    0=\frac{\partial}{\partial s}\frac{\partial^2\bar c}{\partial t^2}
    =R\!\left(\frac{\partial c}{\partial s},\frac{\partial c}{\partial t}\right)
    \frac{\partial c}{\partial t}+\frac{\partial^2}{\partial t^2}\frac{\partial c}{\partial s},
  \]
  so the variational field $J(t)=\frac{\partial\bar c}{\partial s}(0,t)$ solves
  the linear second-order \emph{Jacobi equation} $\ddot J+R(J,\dot c)\dot c=0$.
  A \emph{Jacobi field} along $c$ is any solution of this equation; given
  $J(0),\dot J(0)$ there is a unique such solution (linearity of a second-order
  ODE). Conversely, every pair $(J(0),\dot J(0))$ with $J(0)=0$ arises from the
  geodesic variation $\bar c(s,t)=\exp_p(t(\dot c(0)+sJ(0)))$: since
  $\bar c(s,0)=p$ for all $s$, $J(0)=\frac{\partial\bar c}{\partial s}(0,0)=0$,
  and
  \[
    \dot J(0)=\frac{\partial^2\bar c}{\partial t\,\partial s}(0,0)
    =\frac{\partial^2\bar c}{\partial s\,\partial t}(0,0)
    =\frac{\partial}{\partial s}\big(\dot c(0)+sJ(0)\big)\Big|_{s=0}.
  \]

  \emph{Formalization note — what \texttt{\textbackslash leanok} covers here.} The marker
  certifies the \emph{definition}: \texttt{jacobiEquation} is $\ddot J+R(J,\dot c)\dot c$
  and \texttt{IsJacobiFieldAlong} is its vanishing. It does \emph{not} certify the
  existence-and-uniqueness theorem this node's exposition states in passing (a Jacobi field
  is determined by $(J(0),\dot J(0))$), which is not formalized. The converse realization of
  $J(0)=0$ data by the geodesic variation $\bar c(s,t)=\exp_p(t(\dot c(0)+sJ(0)))$ \emph{is}
  now formalized, in local form, by \cref{rem:pet-ch6-jacobi-dexp-relation}
  (\texttt{jacobiField\_dexp\_relation}).

  The Lean definition is chart-free and needs no chart hypothesis: $\ddot J$ is
  \texttt{derivAlongCurve} applied twice (that operator returns a field along $c$, and is
  chart-independent by \cref{def:pet-ch6-vector-field-along-curve}), and $R$ is the abstract
  \cref{def:pet-ch3-curvature-tensor}. Note the contrast with
  \cref{lem:pet-ch6-third-partial-commutator}, whose \emph{proof} needs a chart and
  \cref{lem:curvature-chart-bridge}.

  \emph{Naming.} The Lean name is \texttt{IsJacobiFieldAlong}, not
  \texttt{IsJacobiField}: the latter is already taken by \cref{def:pet-ch3-jacobi-field}
  for the \emph{distance-function} Jacobi field $L_{\partial_r}J=0$ of §3.2.4 — a different
  notion, with no geodesic and no second derivative. The suffix follows
  \cref{def:pet-ch6-parallel-field}'s \texttt{IsParallelAlong}.
  \source{petersen:page-0253,petersen:page-0254}
\end{definition}

\begin{remark}[Jacobi fields and $D\exp_p$]
  \label{rem:pet-ch6-jacobi-dexp-relation}
  \lean{PetersenLib.jacobiField_dexp_relation}
  \uses{def:pet-ch6-jacobi-field}
  \leanok
  For the geodesic variation $\bar c(s,t)=\exp_p(t(\dot c(0)+sJ(0)))$ of
  \cref{def:pet-ch6-jacobi-field},
  \[
    J(t)=\frac{\partial}{\partial s}\exp_p\big(t(\dot c(0)+sJ(0))\big)\Big|_{s=0}
    =D\exp_p\big(tJ(0)\big),
  \]
  where $tJ(0)\in T_{tv}T_pM$ for $v=\dot c(0)$. Consequently, $D\exp_p$ is
  nonsingular at $t_0v$ if and only if for every $J(t_0)\in
  T_{\exp_p(t_0v)}M$ there is a Jacobi field $J$ along $t\mapsto\exp_p(tv)$
  with $J(0)=0$ and value $J(t_0)$ at $t=t_0$.

  \emph{Formalization note.} The Lean statement is the local form Petersen's
  remark presumes throughout: it produces a normal radius $\rho$ and, for every
  $v$ with $\lvert v\rvert<\rho$ and every $w=J(0)$, a time window $(-\beta,\beta)$
  on which the variation field $\partial_s\bar c(\cdot,s)|_{s=0}$ of
  $\bar c(s,t)=\exp_p(t(v+sw))$ solves the Jacobi equation along
  $t\mapsto\exp_p(tv)$ and vanishes at $t=0$. This is the substantive content: the
  displayed identity $J=D\exp_p(t\,\cdot)$ merely names the variation field, since
  $\partial_s\exp_p(t(v+sw))|_{s=0}=D\exp_p|_{tv}(tw)$ by the chain rule (an
  equality the Lean does not separately restate, as it never mentions $D\exp_p$).
  The exponential rays $t\mapsto\exp_p(t(v+sw))$ are geodesics only while they stay
  in the normal ball, which is exactly what bounds the window. The final
  nonsingularity biconditional ($D\exp_p$ nonsingular at $t_0v$ iff every value
  $J(t_0)$ is realized) is the standard consequence for conjugate points and is not
  separately mechanized here.
  \source{petersen:page-0254,petersen:page-0255}
\end{remark}

\begin{remark}[Jacobi fields and the Hessian of $r$]
  \label{rem:pet-ch6-jacobi-hessian-r}
  \lean{PetersenLib.jacobiField_hess_r}
  \leanok
  \uses{def:pet-ch6-jacobi-field}
  Let $c(t)$ be a unit speed geodesic with $c(0)=p$ and $J$ a Jacobi field
  along $c$ with $J(0)=0$; write $r(x)=|xp|$. As long as $tc'(0)$ lies in the
  interior of the domain on which $r$ is smooth (so that $\dot c(t)=\nabla
  r|_{c(t)}$), realizing $J$ via a geodesic variation $\bar c(s,t)$ as in
  \cref{def:pet-ch6-jacobi-field} gives
  \[
    \operatorname{Hess}r(J(t),J(t))=g(\nabla_{J(t)}\nabla r,J(t))
    =g\!\left(\frac{\partial^2\bar c}{\partial t\,\partial s},J\right)\Big|_{(0,t)}
    =g(\dot J(t),J(t)).
  \]
  The middle equality $g(\nabla_{J}\nabla r,J)=g(\dot J,J)$ is the crux: the
  Lean statement takes $J$ in the \S3.2.4 distance-function sense
  ($L_{\nabla r}J=0$), exactly the condition that turns
  $\nabla_{J}\nabla r$ into $\nabla_{\nabla r}J=\dot J$, and carries Petersen's
  radial hypothesis $\dot c(t)=\nabla r|_{c(t)}$ explicitly. The equation
  $\dot J=\nabla_{\dot c}J$ is the reusable along-curve $\leftrightarrow$ global
  Levi-Civita bridge \lean{PetersenLib.derivAlongCurve_eq_leviCivita_cov}, which
  identifies $\dot J(t)$ (the moving-foot covariant derivative
  \cref{def:pet-ch6-vector-field-along-curve}) with $\nabla_{\dot c(t)}J$ for a
  globally smooth field $J$; the first equality is the Ch.~3
  distance-function computation. The middle geometric interpretation via
  $\partial^2\bar c/\partial t\,\partial s$ is not separately restated, matching
  the honest scoping of \cref{rem:pet-ch6-jacobi-dexp-relation}.
  \source{petersen:page-0255}
\end{remark}

\begin{lemma}[the second variation integrand]
  \label{lem:pet-ch6-second-variation-integrand}
  \lean{PetersenLib.hasDerivAt_chartPairing_transversalAccel_of_geodesic_curvatureTensorAt,PetersenLib.hasDerivAt_chartPairing_transversalAccel_of_geodesic,PetersenLib.covariantDerivCoord_mixedPartial_fst_eq_sub_chartCurvature,PetersenLib.hasDerivAt_chartPairing_slice_ts,PetersenLib.surfaceCovariantDerivS_fderivSnd_eq_surfaceCovariantDerivT_fderivFst}
  \leanok
  \uses{def:pet-ch6-vector-field-along-curve}
  Let $\bar c(s,t)$ be a smooth variation whose base curve $t\mapsto\bar c(0,t)$ is
  a geodesic, and write $V=\partial c/\partial s$ and $T=\partial c/\partial t$.
  Then at $s=0$,
  \[
    \frac{\partial}{\partial t}g\Big(\frac{DV}{\partial s},T\Big)
      =g\Big(\frac{D}{\partial s}\frac{DT}{\partial s},T\Big)-g\big(R(V,T)V,T\big).
  \]
\end{lemma}
\begin{proof}
  \uses{lem:pet-ch6-third-partial-commutator}
  Metric compatibility of the connection along the curve $t\mapsto\bar c(0,t)$
  (\cref{def:pet-ch6-vector-field-along-curve}) gives
  \[
    \frac{\partial}{\partial t}g\Big(\frac{DV}{\partial s},T\Big)
      =g\Big(\frac{D}{\partial t}\frac{DV}{\partial s},T\Big)
       +g\Big(\frac{DV}{\partial s},\frac{DT}{\partial t}\Big).
  \]
  The base curve is a geodesic, so $\frac{DT}{\partial t}=0$ at $s=0$ and the
  second term vanishes.

  For the first, symmetry of the second partials gives
  $\frac{DT}{\partial s}=\frac{DV}{\partial t}$, whence
  $\frac{D}{\partial s}\frac{DT}{\partial s}=\frac{D}{\partial s}\frac{DV}{\partial t}$,
  and \cref{lem:pet-ch6-third-partial-commutator}, applied to the field $V$ in the
  parameters $(s,t)$, gives
  \[
    \frac{D}{\partial s}\frac{DV}{\partial t}-\frac{D}{\partial t}\frac{DV}{\partial s}
      =R(V,T)V.
  \]
  Therefore
  $g\big(\frac{D}{\partial t}\frac{DV}{\partial s},T\big)
    =g\big(\frac{D}{\partial s}\frac{DT}{\partial s},T\big)-g(R(V,T)V,T)$,
  which is the claim.
\end{proof}


\begin{lemma}[the coordinate curvature in one chart, at every point of it]
  \label{lem:pet-ch6-offdiagonal-curvature-bridge}
  \lean{PetersenLib.curvatureTensorAt_chartBasis_of_mem,PetersenLib.curvatureTensor_coordinates_of_mem,PetersenLib.chartCurvatureContraction2_eq_neg_curvatureTensorAt_of_mem,PetersenLib.chartCurvature_eq_curvatureTensorAt_of_mem,PetersenLib.chartCurvature_coordChange}
  \leanok
  \uses{prop:pet-ch3-curvature-coordinate-formula}
  Fix a chart with domain $U$ and coordinate frame $\partial_1,\dots,\partial_n$.
  For \emph{every} $q\in U$ — not merely for one distinguished point of $U$ — the
  curvature tensor of the Levi-Civita connection satisfies
  \[
    R(\partial_i,\partial_j)\partial_k\big|_q=R^l_{ijk}(q)\,\partial_l\big|_q,
    \qquad
    R^l_{ijk}=\partial_i\Gamma^l_{jk}-\partial_j\Gamma^l_{ik}
      +\Gamma^s_{jk}\Gamma^l_{is}-\Gamma^s_{ik}\Gamma^l_{js},
  \]
  where $\Gamma^m_{ij}$ are the Christoffel symbols \emph{of that one chart},
  regarded as functions on $U$ and evaluated at $q$.

  Consequently the coordinate curvature of a chart computes the tensor $R$ at every
  point of the chart's domain, and is natural: if $q$ lies in the domains of two
  charts, their coordinate curvatures at $q$ agree after the change-of-frame
  isomorphism.
  \source{petersen:page-0107}
\end{lemma}
\begin{proof}
  \uses{prop:pet-ch3-curvature-coordinate-formula}
  \leanok
  This is \cref{prop:pet-ch3-curvature-coordinate-formula} read with care about
  where each ingredient is evaluated. That computation uses only two facts, and
  both hold at every $q\in U$ rather than at one point:
  $[\partial_i,\partial_j]=0$, so the bracket term of $R$ drops; and
  $\nabla_{\partial_i}\partial_j=\Gamma^m_{ij}\partial_m$ as an identity of vector
  fields \emph{on $U$}, which is what allows $\nabla_{\partial_i}$ to differentiate
  the coefficient functions $\Gamma^s_{jk}$ by the Leibniz rule. Expanding
  $R(\partial_i,\partial_j)\partial_k=\nabla_{\partial_i}(\Gamma^s_{jk}\partial_s)
  -\nabla_{\partial_j}(\Gamma^s_{ik}\partial_s)$ at $q$ and applying Leibniz gives
  the four displayed terms, each evaluated at $q$.

  The only tempting error is to express the quadratic terms using Christoffel
  symbols \emph{at a distinguished point} of $U$; they must be the coefficient
  functions evaluated at $q$. In particular one may not choose coordinates normal at
  $q$ to kill them, since $q$ ranges over all of $U$ while the chart stays fixed.

  Naturality follows: both charts' coordinate curvatures at $q$ compute the same
  frame-independent tensor $R|_q$, so they agree after the change-of-frame map.
\end{proof}

\begin{theorem}[Theorem 6.1.4 — Synge's second variation formula, 1926]
  \label{thm:pet-ch6-synge-second-variation}
  \leanok
  \lean{PetersenLib.secondVariationEnergy,PetersenLib.secondVariationEnergy_deriv,PetersenLib.intervalIntegrable_syngeIntegrand_window,PetersenLib.syngeIntegrand_eq_chartCurvature,PetersenLib.hasDerivAt_deriv_windowEnergy,PetersenLib.hasDerivAt_deriv_windowEnergy_chart_curvatureTensorAt,PetersenLib.derivAlongCurve_variationField_eq_transfer,PetersenLib.energyFunctional_eventuallyEq_windowEnergy_chart,PetersenLib.secondVariationEnergy_chart_curvatureTensorAt,PetersenLib.secondVariationEnergy_chart,PetersenLib.hasDerivAt_chartPairing_transversalAccel_of_geodesic_curvatureTensorAt_of_mem,PetersenLib.iteratedDeriv_two_windowEnergy_chart,PetersenLib.hasDerivAt_deriv_windowEnergy_chart,PetersenLib.hasDerivAt_integral_chartPairing_ss,PetersenLib.hasDerivAt_chartPairing_slice_ss,PetersenLib.hasDerivAt_windowEnergy_chart_preByParts_shift}
  \uses{def:pet-ch6-vector-field-along-curve,lem:pet-ch6-offdiagonal-curvature-bridge}
  If $\bar c:(-\varepsilon,\varepsilon)\times[a,b]\to M$ is a smooth variation
  of a geodesic $c(t)=\bar c(0,t)$, then
  \[
    \frac{d^2E(c_s)}{ds^2}\bigg|_{s=0}
    =\int_a^b\Big|\frac{\partial^2c}{\partial t\,\partial s}\Big|^2dt
    -\int_a^bg\!\left(R\!\left(\frac{\partial c}{\partial s},\frac{\partial c}{\partial t}\right)
    \frac{\partial c}{\partial t},\frac{\partial c}{\partial s}\right)dt
    +g\!\left(\frac{\partial^2c}{\partial s^2},\frac{\partial c}{\partial t}\right)\bigg|_a^b.
  \]
  \source{petersen:page-0255}
\end{theorem}
\begin{proof}
  \uses{lem:pet-ch6-third-partial-commutator,lem:pet-ch6-second-variation-integrand}
  \leanok
  \emph{From one chart to $[a,b]$: the chart cover.} The identity is first proved for a
  variation contained in a \emph{single} chart, as \texttt{hasDerivAt\_deriv\_windowEnergy}:
  there the chart occurs only as a hypothesis (the slab lies in one chart's source), while the
  conclusion is exactly the display above, written with \texttt{energyFunctional},
  \texttt{transversalAccel}, \texttt{curveVelocity}, \texttt{derivAlongCurve} and
  \texttt{curvatureTensorAt} — no coordinates. That asymmetry is the design: because the chart
  is absent from the conclusion, windows with different chart centres glue.

  Since $\bar c$ is defined on a compact $[a,b]$ it leaves every single chart, so we cover.
  Continuity of $\bar c$ at $(0,\tau)$ gives, at each $\tau\in[a,b]$, a product window that
  $\bar c$ maps into the source of the chart at $c(\tau)$; the $\tau$-balls of these windows
  cover the compact $[a,b]$, so the Lebesgue number lemma yields a single mesh $r$, and a
  uniform partition $a=\tau_0<\cdots<\tau_N=b$ finer than $r/2$ cuts $[a,b]$ into pieces each
  of which — after padding by the mesh, which the $r/2$ choice affords — lies in one chart
  window with an \emph{open} time interval, the shape the window theorem consumes. Shrinking
  the $s$-half-width to the minimum over the finitely many pieces makes one slab serve them all.

  The pieces then telescope. The boundary terms
  $g\big(\frac{\partial^2c}{\partial s^2},\frac{\partial c}{\partial t}\big)\big|_{\tau_j}$
  depend only on $(g,\bar c,t)$ — on neither the window nor the chart — so at an interior
  junction the two abutting pieces contribute the same term with opposite signs and the sum
  collapses to the two outer endpoints; the integrals add by additivity over adjacent
  intervals, using integrability of the Synge integrand on each piece
  (\texttt{intervalIntegrable\_syngeIntegrand\_window}, which reads the integrand in the piece's
  chart via \texttt{syngeIntegrand\_eq\_chartCurvature} and observes it is continuous there).

  One point of care: the argument never forms
  $\operatorname{deriv}\circ\operatorname{deriv}$, because per-piece \emph{equations} about
  $\frac{d^2E_j}{ds^2}\big|_0$ say nothing about their sum. It stays in \texttt{HasDerivAt}
  shape throughout: each piece energy is differentiable at every $s$ near $0$, so
  $\frac{dE}{ds}$ and $\sum_j\frac{dE_j}{ds}$ agree on a \emph{neighbourhood} of $0$ and the
  summed statement transports back along that germ. The book's
  $\frac{d^2E(c_s)}{ds^2}\big|_{s=0}=\dots$ shape is then recovered as a corollary
  (\texttt{secondVariationEnergy\_deriv}).

  \emph{Route.} Petersen differentiates the first variation \emph{after} its
  integration by parts. The following variant instead differentiates it
  \emph{before}, and is the one formalized. Both give the identity below; the
  second is shorter and, unlike the first, never differentiates a boundary term.
  Write $T=\partial c/\partial t$, $V=\partial c/\partial s$. Differentiating
  $E(c_s)=\tfrac12\int_a^b g(T,T)\,dt$ under the integral once gives, for every
  $s$,
  \[
    \frac{dE(c_s)}{ds}=\int_a^b g\Big(\frac{D T}{\partial s},T\Big)dt,
  \]
  a \emph{pure integral, with no boundary term}. Differentiating once more and
  using metric compatibility,
  \[
    \frac{d^2E(c_s)}{ds^2}
    =\int_a^b g\Big(\frac{D}{\partial s}\frac{D T}{\partial s},T\Big)dt
     +\int_a^b\Big|\frac{D T}{\partial s}\Big|^2dt .
  \]
  Now $\frac{DT}{\partial s}=\frac{DV}{\partial t}$ by symmetry of second
  partials, so the second integrand is $|\dot V|^2$. For the first, apply
  \cref{lem:pet-ch6-third-partial-commutator} to $V$:
  $\frac{D}{\partial s}\frac{D V}{\partial t}
   =\frac{D}{\partial t}\frac{D V}{\partial s}+R(V,T)V$. At $s=0$ the curve is a
  geodesic, $\frac{DT}{\partial t}=0$, whence
  $g\big(\frac{D}{\partial t}\frac{DV}{\partial s},T\big)
   =\frac{\partial}{\partial t}g\big(\frac{DV}{\partial s},T\big)$, which the
  fundamental theorem of calculus turns into the boundary term
  $g\big(\frac{\partial^2c}{\partial s^2},T\big)\big|_a^b$. Finally
  $g(R(V,T)V,T)=-g(R(V,T)T,V)$ by antisymmetry of $R$ in its last two slots,
  giving the curvature integral. This is exactly the claimed formula.

  \emph{Petersen's own argument.}
  By the first variation formula,
  \[
    \frac{dE(c_s)}{ds}=-\int_a^bg\!\left(\frac{\partial c}{\partial s},
    \frac{\partial^2c}{\partial t^2}\right)dt
    +g\!\left(\frac{\partial c}{\partial s},\frac{\partial c}{\partial t}\right)\Big|_{(s,a)}^{(s,b)}.
  \]
  Differentiating in $s$ and using the product rule from
  \cref{def:pet-ch6-vector-field-along-curve},
  \begin{align*}
    \frac{\partial^2E(c_s)}{\partial s^2}
    &=-\int_a^bg\!\left(\frac{\partial^2c}{\partial s^2},\frac{\partial^2c}{\partial t^2}\right)dt
    -\int_a^bg\!\left(\frac{\partial c}{\partial s},
    \frac{\partial^3c}{\partial s\,\partial t^2}\right)dt\\
    &\quad+g\!\left(\frac{\partial^2c}{\partial s^2},\frac{\partial c}{\partial t}\right)\Big|_{(s,a)}^{(s,b)}
    +g\!\left(\frac{\partial c}{\partial s},\frac{\partial^2c}{\partial s\,\partial t}\right)\Big|_{(s,a)}^{(s,b)}.
  \end{align*}
  Setting $s=0$, using that $c(0,t)$ is a geodesic (so
  $\frac{\partial^2c}{\partial t^2}=0$ at $s=0$), and applying
  \cref{lem:pet-ch6-third-partial-commutator} to rewrite
  $\frac{\partial^3c}{\partial s\,\partial t^2}
  =\frac{\partial^3c}{\partial t\,\partial s\,\partial t}
  -R\!\left(\frac{\partial c}{\partial s},\frac{\partial c}{\partial t}\right)\frac{\partial c}{\partial t}$
  gives
  \begin{align*}
    \frac{\partial^2E(c_s)}{\partial s^2}\bigg|_{s=0}
    &=-\int_a^bg\!\left(\frac{\partial c}{\partial s},R\!\left(\frac{\partial c}{\partial s},
    \frac{\partial c}{\partial t}\right)\frac{\partial c}{\partial t}\right)dt
    +\int_a^bg\!\left(\frac{\partial^2c}{\partial t\,\partial s},\frac{\partial^2c}{\partial t\,\partial s}\right)dt\\
    &\quad-\int_a^b\frac{\partial}{\partial t}g\!\left(\frac{\partial c}{\partial s},
    \frac{\partial^2c}{\partial s\,\partial t}\right)dt
    +g\!\left(\frac{\partial^2c}{\partial s^2},\frac{\partial c}{\partial t}\right)\Big|_a^b
    +g\!\left(\frac{\partial c}{\partial s},\frac{\partial^2c}{\partial s\,\partial t}\right)\Big|_a^b.
  \end{align*}
  The fundamental theorem of calculus cancels the last integral against the
  last boundary term, leaving exactly the claimed formula.
  \source{petersen:page-0255,petersen:page-0256}
\end{proof}

\begin{remark}[special cases of the second variation formula]
  \label{rem:pet-ch6-second-variation-special-cases}
  \leanok
  \lean{PetersenLib.secondVariationEnergy_properVariation,PetersenLib.secondVariationEnergy_parallelField,PetersenLib.secondVariationEnergy_parallelField_geodesicTransversals,PetersenLib.curveAcceleration_const}
  \uses{thm:pet-ch6-synge-second-variation,def:pet-ch6-parallel-field}
  For a \emph{proper} variation (fixed endpoints) the boundary term in
  \cref{thm:pet-ch6-synge-second-variation} drops out, and with $V=\frac{\partial
  c}{\partial s}(0,\cdot)$,
  \[
    \frac{d^2E(c_s)}{ds^2}\bigg|_{s=0}=\int_a^b|\dot V|^2dt-\int_a^bg(R(V,\dot c)\dot c,V)dt.
  \]
  If instead $V$ is parallel, the first integral drops out:
  \[
    \frac{d^2E(c_s)}{ds^2}\bigg|_{s=0}=-\int_a^bg(R(V,\dot c)\dot c,V)dt
    +g\!\left(\frac{\partial^2c}{\partial s^2},\dot c\right)\bigg|_a^b,
  \]
  and if in addition $s\mapsto\bar c(s,t)$ is a geodesic for each $t$, the
  remaining boundary term also vanishes.
  \source{petersen:page-0256}
\end{remark}

\section{Nonpositive Sectional Curvature}

\begin{lemma}[Lemma 6.2.1]
  \label{lem:pet-ch6-nonsingular-exp-covering}
  \lean{PetersenLib.exp_nonsingular_isCoveringMap}
  If $\exp_p:T_pM\to M$ is nonsingular everywhere (no critical points), then it
  is a covering map.
  \source{petersen:page-0257}
\end{lemma}
\begin{proof}
  By definition $\exp_p$ is then an immersion, so pulling back the metric on
  $M$ makes $\exp_p$ a local Riemannian isometry on $T_pM$. A local isometry
  from a complete Riemannian manifold is a covering map, so it suffices that
  the pullback metric on $T_pM$ is complete; this holds because it is
  geodesically complete at the origin, straight lines through the origin
  remaining geodesics for the pullback metric.
  \source{petersen:page-0257}
\end{proof}

\begin{theorem}[Theorem 6.2.2 — Mangoldt 1881, Hadamard 1889, Cartan 1925]
  \label{thm:pet-ch6-cartan-hadamard}
  \lean{PetersenLib.cartanHadamard}
  \uses{lem:pet-ch6-nonsingular-exp-covering,def:pet-ch6-sec-bounds}
  If $(M,g)$ is complete, connected, and has $\sec\le0$, then the universal
  cover of $M$ is diffeomorphic to $\mathbb{R}^n$.
  \source{petersen:page-0257}
\end{theorem}
\begin{proof}
  \uses{def:pet-ch6-jacobi-field, rem:pet-ch6-jacobi-dexp-relation, lem:pet-ch6-nonsingular-exp-covering}
  It suffices to show $|D\exp_p(w)|>0$ for every nonzero $w\in T_pT_pM$;
  \cref{lem:pet-ch6-nonsingular-exp-covering} then shows $\exp_p$ is a
  covering map, and $T_pM\cong\mathbb{R}^n$ is simply connected, so it is the
  universal cover. Fix $p$, $w\ne0$, and (using
  \cref{rem:pet-ch6-jacobi-dexp-relation}) the Jacobi field $J$ along
  $c(t)=\exp_p(tw)$ with $J(0)=0$, $\dot J(0)=w$, so $|D\exp_p(w)|=|J(1)|$.
  Then
  \[
    \frac{d}{dt}\Big(\tfrac12|J(t)|^2\Big)=g(\dot J,J),\qquad
    \frac{d^2}{dt^2}\Big(\tfrac12|J(t)|^2\Big)
    =g(\ddot J,J)+|\dot J|^2=-g(R(J,\dot c)\dot c,J)+|\dot J|^2\ge|\dot J|^2,
  \]
  the last step using $\sec\le0$, i.e.\ $g(R(x,y)y,x)\le0$ for all $x,y$.
  Integrating from $0$ (where $g(\dot J(0),J(0))=0$),
  $g(\dot J(t),J(t))\ge\int_0^t|\dot J|^2\,dt\ge0$, with equality for all $t$
  forcing $\dot J\equiv0$ hence $w=\dot J(0)=0$, contrary to hypothesis. So for
  $w\ne0$ this derivative is $>0$ on $(0,1]$ (it cannot vanish identically),
  and integrating once more, $\tfrac12|J(t)|^2>0$ for $t\in(0,1]$; in
  particular $|J(1)|>0$, as required.
  \source{petersen:page-0257,petersen:page-0258}
\end{proof}

\begin{remark}[failure for Ricci or scalar curvature bounds]
  \label{rem:pet-ch6-cartan-hadamard-fails-ricci}
  \lean{PetersenLib.cartanHadamard_fails_forRicci}
  \uses{thm:pet-ch6-cartan-hadamard}
  No analogue of \cref{thm:pet-ch6-cartan-hadamard} holds under $\operatorname{Ric}\le0$
  or $\operatorname{scal}\le0$ in place of $\sec\le0$: there exist Ricci-flat metrics on
  $\mathbb{R}^2\times S^{n-2}$ and scalar-flat metrics on $\mathbb{R}\times S^{n-1}$.
  \source{petersen:page-0258}
\end{remark}

\begin{definition}[strictly convex function]
  \label{def:pet-ch6-convex-function}
  \lean{PetersenLib.IsConvexOn,PetersenLib.IsStrictlyConvexOn}
  \leanok
  A function on a Riemannian manifold is \emph{(strictly) convex} if its
  restriction to every \emph{nonconstant} geodesic is (strictly) convex as a
  function of the affine parameter. The Lean statements are relative to an open
  set $U$ (``convex on $U$''), quantifying over the nonconstant geodesics that
  stay in $U$; this is what \Cref{thm:pet-ch6-convexity-radius-criterion}
  requires, and the global notion is the case $U=M$.

  \emph{Correction.} Petersen says ``every geodesic''. Taken literally that
  makes \emph{strict} convexity vacuous --- unsatisfiable by \emph{every}
  function --- because a constant curve is a geodesic ($\ddot c=0$), so
  $f\circ c$ is constant and hence not strictly convex, for any $f$ whatsoever.
  Restricting to nonconstant geodesics is Petersen's own intent, as his
  argument in \Cref{rem:pet-ch6-hessian-f0-positive-definite} shows: he derives
  strict convexity from $(f_0\circ c)''=\operatorname{Hess}f_0(\dot c,\dot c)>0$,
  which needs $\dot c\neq0$. The Lean carries the guard as
  \texttt{IsNonconstantGeodesicOn}.
  \source{petersen:page-0259}
\end{definition}

\begin{remark}[$\operatorname{Hess}f_0\succ0$ in nonpositive curvature]
  \label{rem:pet-ch6-hessian-f0-positive-definite}
  \lean{PetersenLib.hess_f0_posDef_nonpositiveCurvature}
  \leanok
  \uses{rem:pet-ch6-jacobi-hessian-r,def:pet-ch6-convex-function,thm:pet-ch6-cartan-hadamard}
  Let $(M,g)$ be complete, simply connected, and of nonpositive curvature, so
  by \cref{thm:pet-ch6-cartan-hadamard} the distance $r(x)=|xp|$ is smooth on
  $M-\{p\}$ and $f_0(x)=\tfrac12r(x)^2$ is smooth everywhere, with
  $\operatorname{Hess}f_0=dr^2+r\operatorname{Hess}r$. If $J$ is a Jacobi field along a unit
  speed geodesic from $p$ with $J(0)=0$, then by
  \cref{rem:pet-ch6-jacobi-hessian-r} and the estimate in the proof of
  \cref{thm:pet-ch6-cartan-hadamard},
  \[
    \operatorname{Hess}r(J(b),J(b))=g(\dot J(b),J(b))\ge\int_0^b|\dot J|^2dt>0
  \]
  for $b>0$ and $J(b)\ne0$ (any $J(b)$ arises this way), so $\operatorname{Hess}r$, hence
  $\operatorname{Hess}f_0$, is positive definite. Consequently $f_0\circ c$ is (strictly)
  convex along any geodesic $c$, since
  $\frac{d^2}{dt^2}f_0\circ c=\operatorname{Hess}f_0(\dot c,\dot c)>0$; thus $f_0$ is a
  strictly convex function in the sense of
  \cref{def:pet-ch6-convex-function}.

  \emph{Formalization note.} The \texttt{\textbackslash leanok} content
  (\texttt{hess\_f0\_posDef\_nonpositiveCurvature}) is the genuinely-new analytic
  core: the \emph{strict positivity} $0<g(\dot J(b),J(b))$ of the displayed chain,
  for every along-curve Jacobi field ($\texttt{IsJacobiFieldAlong}$) with $J(0)=0$,
  $J(b)\ne0$, $b>0$, under $\sec\le0$ ($\texttt{HasSecBoundedAbove}\ g.\texttt{leviCivita}\ 0$).
  Via the already-\texttt{\textbackslash leanok} sibling
  \cref{rem:pet-ch6-jacobi-hessian-r} ($\operatorname{Hess}r(J,J)=g(\dot J,J)$) this
  reads $\operatorname{Hess}r(J(b),J(b))>0$. The proof runs entirely in the
  along-curve language: with $\varphi(t)=g(J,\dot J)$ one has
  $\varphi'=g(\dot J,\dot J)-g(R(J,\dot c)\dot c,J)\ge0$ (metric nonnegativity and
  $\sec\le0$, the denominator-cleared bound
  $\texttt{curvatureTensorFourAt\_le}$), so $\varphi$ is monotone; the fundamental
  theorem of calculus applied to $g(J,J)$ gives $\int_0^b2\varphi=|J(b)|^2>0$, which
  forces $\varphi(b)>0$. The middle inequality
  $g(\dot J(b),J(b))\ge\int_0^b|\dot J|^2$ — ``the estimate in the proof of
  \cref{thm:pet-ch6-cartan-hadamard}'' — is the separate lemma
  \texttt{jacobiField\_energy\_le\_boundary\_nonpositiveCurvature} (from
  \cref{rem:pet-ch6-ex-6-7-25}'s boundary identity plus $-g(R(J,\dot c)\dot c,J)\ge0$).
  The two remaining links to full positive-definiteness of $\operatorname{Hess}f_0$
  ---the surjectivity ``any $J(b)$ arises this way'' (\cref{thm:pet-ch6-cartan-hadamard}'s
  nonsingular $D\exp_p$) and the decomposition $\operatorname{Hess}f_0=dr^2+r\operatorname{Hess}r$---
  are the covering-theory pieces scoped out, exactly as in
  \cref{rem:pet-ch6-jacobi-hessian-r}.
  \source{petersen:page-0259}
\end{remark}

\begin{remark}[maxima of convex functions and unique minima]
  \label{rem:pet-ch6-max-of-convex-functions}
  \lean{PetersenLib.max_isConvex,PetersenLib.strictlyConvex_uniqueMinimum}
  \leanok
  \uses{def:pet-ch6-convex-function,thm:pet-ch5-hopf-rinow}
  The maximum of finitely many convex functions is convex (this reduces to
  the one-dimensional statement by restricting to geodesics). Any proper,
  nonnegative, strictly convex function $f$ on a complete manifold has a
  unique minimum: properness and boundedness below give a minimum, and if
  there were two minima, strict convexity restricted to a geodesic joining
  them would force smaller values on the interior of the segment than at
  either endpoint, a contradiction.

  \emph{Formalization notes.} (i) ``Maximum of finitely many'' is
  \texttt{Finset.sup'} over a \emph{nonempty} \texttt{Finset}: a maximum over
  an empty family has no value in $\mathbb{R}$. The $k$-point family of
  \cref{def:pet-ch6-linfty-center-of-mass} is nonempty, so this costs the
  consumer nothing. (ii) The Lean adds two hypotheses Petersen leaves
  implicit: $f$ is \emph{continuous} (properness alone does not give an
  attained infimum; Petersen's $f$ is a max of the smooth $\tfrac12r_{p_i}^2$,
  so this is free at the point of use, and deriving it from convexity would
  need convex~$\Rightarrow$~locally Lipschitz, which this project has not
  formalized), and $M$ is \emph{connected} (Petersen's standing assumption,
  and what \cref{thm:pet-ch5-hopf-rinow} needs to produce the joining
  geodesic). Properness is spelled as compact sublevel sets, which for
  nonnegative $f$ is equivalent to $f$ being a proper map. (iii)
  Nonnegativity is stated for fidelity but is not used --- compact sublevel
  sets already give existence.
  \source{petersen:page-0259}
\end{remark}

\begin{definition}[$L^\infty$ center of mass]
  \label{def:pet-ch6-linfty-center-of-mass}
  \lean{PetersenLib.centerOfMassLinfty}
  \uses{rem:pet-ch6-hessian-f0-positive-definite,rem:pet-ch6-max-of-convex-functions}
  Let $(M,g)$ be complete, simply connected, of nonpositive curvature, and
  $p_1,\dots,p_k\in M$. By
  \cref{rem:pet-ch6-hessian-f0-positive-definite,rem:pet-ch6-max-of-convex-functions}
  the function $x\mapsto\max\{f_{0,p_1}(x),\dots,f_{0,p_k}(x)\}$ is proper,
  nonnegative, and strictly convex, hence has a unique minimum, denoted
  $\operatorname{cm}_\infty\{p_1,\dots,p_k\}$ and called the \emph{$L^\infty$ center of
  mass} of $\{p_1,\dots,p_k\}$: it is the center of the smallest closed ball
  containing $\{p_1,\dots,p_k\}$.
  \source{petersen:page-0259,petersen:page-0260}
\end{definition}

\begin{theorem}[Theorem 6.2.3 — Cartan, 1925]
  \label{thm:pet-ch6-cartan-fixed-point}
  \lean{PetersenLib.cartan_finiteOrderIsometry_fixedPoint}
  \uses{def:pet-ch6-linfty-center-of-mass}
  If $(M,g)$ is a complete simply connected Riemannian manifold of nonpositive
  curvature, then any isometry $F:M\to M$ of finite order has a fixed point.
  \source{petersen:page-0260}
\end{theorem}
\begin{proof}
  Let $k$ be the order of $F$ (the smallest positive integer with $F^k=\mathrm
  {id}$), fix $p\in M$, and set
  $q=\operatorname{cm}_\infty\{p,F(p),\dots,F^{k-1}(p)\}$, the unique minimizer of
  $f(x)=\max\{f_{0,p}(x),\dots,f_{0,F^{k-1}(p)}(x)\}$. Since $F$ is an
  isometry,
  \begin{align*}
    f(F(q))&=\tfrac12\big(\max\{|F(q)p|,\dots,|F(q)F^{k-1}(p)|\}\big)^2\\
    &=\tfrac12\big(\max\{|F(q)F^k(p)|,\dots,|F(q)F^{k-1}(p)|\}\big)^2
    =\tfrac12\big(\max\{|qF^{k-1}(p)|,\dots,|qF^{k-2}(p)|\}\big)^2=f(q),
  \end{align*}
  (reindexing the finite orbit $\{p,\dots,F^{k-1}(p)\}=\{F(p),\dots,F^k(p)\}$).
  So $F(q)$ is also a minimum of the strictly convex function $f$; by
  uniqueness of the minimizer (\cref{def:pet-ch6-linfty-center-of-mass}),
  $F(q)=q$.
  \source{petersen:page-0260}
\end{proof}

\begin{corollary}[Corollary 6.2.4]
  \label{cor:pet-ch6-torsion-free-fundamental-group}
  \lean{PetersenLib.fundamentalGroup_torsionFree_nonpositiveCurvature}
  \uses{thm:pet-ch6-cartan-fixed-point}
  If $(M,g)$ is a complete Riemannian manifold of nonpositive curvature, then
  $\pi_1(M)$ is torsion free: every nontrivial element has infinite order.
  \source{petersen:page-0261}
\end{corollary}
\begin{proof}
  \uses{thm:pet-ch6-cartan-fixed-point}
  A nontrivial finite-order element of $\pi_1(M)$ acts on the universal cover
  $\tilde M$ (complete, simply connected, of the same nonpositive curvature) as
  a nontrivial deck transformation of finite order. By
  \cref{thm:pet-ch6-cartan-fixed-point} this deck transformation has a fixed
  point; but a nontrivial deck transformation of a covering acts freely, a
  contradiction. Hence no nontrivial element of $\pi_1(M)$ has finite order.
  \source{petersen:page-0261}
\end{proof}

\begin{lemma}[Lemma 6.2.5]
  \label{lem:pet-ch6-hessian-lower-bound}
  \lean{PetersenLib.hess_f0_ge_metric_nonpositiveCurvature}
  \uses{rem:pet-ch6-hessian-f0-positive-definite}
  If $(M,g)$ has nonpositive curvature, then any modified distance function
  $f_0=\tfrac12r^2$ satisfies $\operatorname{Hess}f_0\ge g$.
  \source{petersen:page-0261}
\end{lemma}
\begin{proof}
  \uses{rem:pet-ch6-jacobi-hessian-r}
  Writing $g=dr^2+g_r$ and $\operatorname{Hess}f_0=dr^2+r\operatorname{Hess}r$, the claim reduces
  to $r\operatorname{Hess}r\ge g_r$, i.e.\ to
  $t\cdot g(\dot J(t),\dot J(t))=t\cdot\operatorname{Hess}r(J(t),J(t))\ge g(J(t),J(t))$
  for a Jacobi field $J$ along a unit speed geodesic from $p$ with $J(0)=0$,
  $\dot J(0)\perp\dot c(0)$. Consider
  $\lambda(t)=|J(t)|^2/g(\dot J(t),J(t))$ (equivalently
  $t-\lambda(t)\ge0$ is to be shown). By l'H\^opital's rule, using
  $g(\ddot J,J)+|\dot J|^2=-g(R(J,\dot c)\dot c,J)+|\dot J|^2$,
  \[
    \lambda(0)=\frac{2g(\dot J(0),J(0))}{-g(R(J(0),\dot c(0))\dot c(0),J(0))+|\dot J(0)|^2}
    =\frac{0}{|\dot J(0)|^2}=0.
  \]
  Its derivative is
  \[
    \dot\lambda(t)=\frac{2\big(g(J,\dot J)\big)^2-|\dot J|^2|J|^2+g(R(J,\dot c)\dot c,J)|J|^2}
    {\big(g(J,\dot J)\big)^2}.
  \]
  Since $\sec\le0$, $g(R(J,\dot c)\dot c,J)\le0$, so the numerator is at most
  $2(g(J,\dot J))^2-|\dot J|^2|J|^2$; by Cauchy--Schwarz,
  $(g(J,\dot J))^2\le|J|^2|\dot J|^2$, so
  \[
    \dot\lambda(t)\le\frac{2\big(g(J,\dot J)\big)^2-\big(g(J,\dot J)\big)^2}{\big(g(J,\dot J)\big)^2}=1.
  \]
  Integrating from $\lambda(0)=0$ gives $\lambda(t)\le t$, i.e.\
  $t\cdot g(\dot J(t),J(t))\ge|J(t)|^2$, which is the desired inequality; the
  claim about $\operatorname{Hess}f_0$ follows.
  \source{petersen:page-0261,petersen:page-0262}
\end{proof}

\begin{remark}[triangle and quadrilateral comparison]
  \label{rem:pet-ch6-triangle-comparison}
  \lean{PetersenLib.triangleComparison_nonpositiveCurvature}
  \uses{lem:pet-ch6-hessian-lower-bound}
  Integrating $\frac{d}{dt}(f_{0,p}\circ c)\ge1$ along a unit speed geodesic
  $c$ (using \cref{lem:pet-ch6-hessian-lower-bound}) yields, for a triangle in
  $M$ with side lengths $a,b,c$ and angle $\alpha$ opposite $a$,
  \[
    a^2\ge b^2+c^2-2bc\cos\alpha.
  \]
  Consequently the angle sum in any geodesic triangle is $\le\pi$, and in any
  geodesic quadrilateral is $\le2\pi$. If $\sec<0$, these inequalities are
  strict unless the base point $p$ lies on the geodesic $c$; in particular the
  angle sum in any nondegenerate quadrilateral is $<2\pi$.
  \source{petersen:page-0262}
\end{remark}

\begin{definition}[axis, period, displacement function]
  \label{def:pet-ch6-axis-displacement}
  \lean{PetersenLib.IsAxis,PetersenLib.IsAxisPeriod,PetersenLib.axisPeriod,
    PetersenLib.displacementFunction}
  \leanok
  For an isometry $F:M\to M$, an \emph{axis} for $F$ is a geodesic
  $c:\mathbb{R}\to M$ such that $F\circ c$ is a reparametrization of $c$; since
  isometries map geodesics to geodesics, $F\circ c(t)=c(\pm t+a)$ for some
  $a\in\mathbb{R}$. If the sign is $-$, $c(a/2)$ is fixed by $F$. When
  $F\circ c(t)=c(t+a)$, $a$ is the \emph{period} of $F$ with respect to $c$
  (depending on the parametrization). The \emph{displacement function} of $F$
  is $x\mapsto\delta_F(x)=|xF(x)|$.

  Since Petersen's period is explicitly not unique, the relation
  ``$a$ is a period of $F$ with respect to $c$'' is primary and is the
  \texttt{Prop} \texttt{IsAxisPeriod}; \texttt{axisPeriod} is the
  derived $\mathbb{R}$-valued \emph{choice}, the infimum of the
  \emph{positive} periods. The positivity restriction is not cosmetic: over all
  periods the infimum would be junk on every periodic axis, since the period
  set is then $\{a_0+kP:k\in\mathbb{Z}\}$, unbounded below, on which
  $\operatorname{sInf}$ returns $0$ --- reporting ``$F$ fixes $c$ pointwise''
  for exactly the nontrivial translations of \Cref{thm:pet-ch6-preissmann}.
  Petersen's periods are positive in any case: \Cref{lem:pet-ch6-axis-existence}
  produces an axis of period $1$ from a unit-speed segment.
  \source{petersen:page-0263}
\end{definition}

\begin{lemma}[Lemma 6.2.7]
  \label{lem:pet-ch6-axis-existence}
  \lean{PetersenLib.axis_of_positiveMinimalDisplacement}
  \uses{def:pet-ch6-axis-displacement}
  If $F:M\to M$ is an isometry of a complete Riemannian manifold and
  $\delta_F$ has a positive minimum, then $F$ has an axis.
  \source{petersen:page-0263}
\end{lemma}
\begin{proof}
  Let $\delta_F$ have a minimum at $p$ and $c:[0,1]\to M$ a unit-speed segment
  from $p$ to $F(p)$; then $F\circ c$ is a segment of the same speed from
  $F(p)$ to $F^2(p)$. For fixed $t\in[0,1]$,
  \begin{align*}
    \delta_F(p)&\le\delta_F(c(t))=|c(t)(F\circ c)(t)|
    \le|c(t)c(1)|+|c(1)(F\circ c)(t)|\\
    &=|c(t)c(1)|+|(F\circ c)(0)(F\circ c)(t)|
    =|c(t)c(1)|+|c(0)c(t)|=|c(0)c(1)|=|pF(p)|=\delta_F(p),
  \end{align*}
  using that $F$ is an isometry to identify $(F\circ c)(0)=F(p)=c(1)$. Hence
  equality holds throughout, so $c|_{[t,1]}$ followed by $F\circ c|_{[0,t]}$
  realizes the distance $|c(t)(F\circ c)(t)|$ via the concatenated path of
  the same total length as the triangle inequality bound, forcing this
  concatenation to be a segment, hence a geodesic. Thus $c|_{[0,1]}$ and
  $F\circ c|_{[0,1]}$ fit together as the extension of $c$ to $[0,2]$, i.e.\
  $(F\circ c)(t)=c(1+t)$. Repeating forwards and backwards extends $c$ to an
  axis for $F$ on all of $\mathbb{R}$, of period $1$.
  \source{petersen:page-0263}
\end{proof}

\begin{definition}[deck transformation]
  \label{def:pet-ch6-deck-transformation}
  \lean{PetersenLib.DeckTransformation,PetersenLib.DeckTransformation.isRiemannianIsometry}
  \leanok
  For the universal cover $\pi:\tilde M\to M$, a \emph{deck transformation} is
  a map $F:\tilde M\to\tilde M$ with $\pi\circ F=\pi$; it is determined by the
  value $F(p)\in\pi^{-1}(q)$ for a chosen $p\in\pi^{-1}(q)$, and the
  fundamental group $\pi_1(M,q)$ acts by deck transformations: a loop
  $[\alpha]\in\pi_1(M,q)$ and lift $\tilde\alpha$ with
  $\tilde\alpha(0)=p$ produce the deck transformation with
  $F(p)=\tilde\alpha(1)$. In the Riemannian setting, deck transformations are
  isometries of $\tilde M$ since $\pi$ is a local isometry.

  The Lean states this for an \emph{arbitrary} smooth covering $\pi:M\to M'$
  rather than specifically the universal cover --- the universal cover is
  available in neither Mathlib nor this project (see
  \Cref{thm:pet-ch6-cartan-hadamard}) --- and takes $F$ to be a diffeomorphism,
  which is automatic for deck transformations of a smooth covering. The
  operative closing claim, that deck transformations are isometries, is the
  proved \texttt{DeckTransformation.isRiemannianIsometry}, for the metric
  pulled back along $\pi$; it is the hinge \Cref{lem:pet-ch6-deck-transformation-dilation}
  needs. The remaining prose --- the $\pi_1(M,q)$-action and the ``determined by
  $F(p)$'' rigidity --- is \emph{not} formalized: both need covering-space
  lifting/monodromy machinery rather than Riemannian geometry.
  \source{petersen:page-0264}
\end{definition}

\begin{lemma}[Lemma 6.2.8]
  \label{lem:pet-ch6-deck-transformation-dilation}
  \lean{PetersenLib.deckTransformation_minimalDisplacement_closedGeodesic}
  \uses{def:pet-ch6-deck-transformation,def:pet-ch6-axis-displacement,lem:pet-ch6-axis-existence}
  If $F:\tilde M\to\tilde M$ is a nontrivial deck transformation on the
  universal cover over a compact base $M$, then $\delta_F$ has a positive
  minimum; the axis through the minimizing point projects to a closed
  geodesic in $M$ of minimal length in its free homotopy class, and
  $\delta_F(x)\ge2\operatorname{inj}_{\pi(x)}(M)$.
  \source{petersen:page-0264}
\end{lemma}
\begin{proof}
  \emph{Loops generated by $F$.} If $x_i\in\tilde M$ ($i=0,1$) are joined to
  $F(x_i)$ by curves $c_i:[0,1]\to\tilde M$, the loops $\pi\circ c_i$ are
  freely homotopic in $M$: extend $H(s,0):[0,1]\to M$ from $x_0$ to $x_1$,
  set $H(s,1)=F(H(s,0))$ (lifted arbitrarily) and $H(i,t)=c_i(t)$; simple
  connectivity of $\tilde M$ extends this boundary data on $\partial[0,1]^2$
  to a map $H:[0,1]^2\to\tilde M$, and $\pi\circ H$ is the required homotopy
  since $\pi(H(s,1))=\pi(F(H(s,0)))=\pi(H(s,0))$. Conversely, if a loop at
  $\pi(x_1)$ freely homotopic to $\pi\circ c_0$ lifts (via the homotopy
  $\tilde H$ with $\tilde H(0,\cdot)=c_0$) then, since both $\tilde H(s,0)$
  and $\tilde H(s,1)$ lift $H(s,0)$ and $F\circ\tilde H(\cdot,0)$ agrees with
  $\tilde H(\cdot,1)$ at $s=0$, uniqueness of lifts gives
  $F\circ\tilde H(s,0)=\tilde H(s,1)$; at $s=1$ this shows the loop lifts to a
  curve from $x_1$ to $F(x_1)$.

  \emph{No small loops.} Since $F$ is nontrivial, none of these loops is
  homotopically trivial; hence $\delta_F(x)\ge2\operatorname{inj}_{\pi(x)}(M)$, as a
  shorter segment from $x$ to $F(x)$ would generate a loop of length
  $<2\operatorname{inj}_{\pi(x)}(M)$ contained in and hence contractible within
  $B(\pi(x),\operatorname{inj}_{\pi(x)}(M))$.

  \emph{Minimizing the displacement.} Take $q_i\in\tilde M$ with
  $\delta_F(q_i)\to\inf\delta_F\ge\operatorname{inj}M$ and segments
  $\tilde c_i:[0,1]\to\tilde M$ from $q_i$ to $F(q_i)$; let
  $c_i=\pi\circ\tilde c_i$. Compactness of $M$ and $|c_i|=\delta_F(q_i)$ allow
  passing to a subsequence with $\dot c_i(0)\to v\in T_qM$, $q=\lim c_i(0)$,
  $|v|=\inf\delta_F$, so $c_i\to c(t)=\exp_q(tv)$ by continuity of $\exp$.
  For large $i$, $|c_i(t)c(t)|<\operatorname{inj}(M)$, so $c_i$ and $c$ are freely
  homotopic via the unique short geodesics joining corresponding points; by
  the characterization above, the lift $\tilde c$ of $c$ satisfies
  $F(\tilde c(0))=\tilde c(1)$, and
  $\delta_F(\tilde c(0))\le L(\tilde c)=L(c)=|v|=\inf\delta_F$, so equality
  holds and $c$ has minimal length in its free homotopy class. The first
  variation formula then shows $c$ is a closed geodesic.
  \source{petersen:page-0264,petersen:page-0265}
\end{proof}

\begin{theorem}[Theorem 6.2.6 — Preissmann, 1943]
  \label{thm:pet-ch6-preissmann}
  \lean{PetersenLib.preissmannTheorem}
  \uses{lem:pet-ch6-deck-transformation-dilation,def:pet-ch6-axis-displacement,rem:pet-ch6-triangle-comparison}
  If $(M,g)$ is a compact manifold of negative curvature, then every Abelian
  subgroup of $\pi_1(M)$ is cyclic. In particular, no compact product manifold
  $M\times N$ admits a metric of negative curvature.
  \source{petersen:page-0263}
\end{theorem}
\begin{proof}
  \uses{lem:pet-ch6-deck-transformation-dilation, rem:pet-ch6-triangle-comparison}
  By \cref{lem:pet-ch6-deck-transformation-dilation} every nontrivial deck
  transformation $F$ on the universal cover $\tilde M$ has an axis. In
  negative curvature axes are unique: if $c_1\ne c_2$ were both axes for $F$
  and intersected once, $F$-invariance would force a second intersection and
  hence coincidence of the geodesics; so distinct axes are disjoint. Joining
  $p_1\in c_1$, $p_2\in c_2$ by a segment $\sigma$, the quadrilateral
  $p_1,p_2,F(p_1),F(p_2)$ (with $F\circ\sigma$ the fourth side, since $F$
  preserves $c_1,c_2$) has adjacent angles along the two axes summing to
  $\pi$ each, hence total angle sum $2\pi$; by
  \cref{rem:pet-ch6-triangle-comparison} this forces the quadrilateral to be
  degenerate, i.e.\ all four points collinear — contradicting $c_1\ne c_2$.
  So every nontrivial deck transformation has a unique axis.

  Now let $G$ commute with $F$, and let $c$ be the axis of $F$ with period
  $l$: $(G\circ c)(t+l)=(G\circ F\circ c)(t)=(F\circ G\circ c)(t)$ shows
  $G\circ c$ is again an axis for $F$, hence (by uniqueness) $G\circ c=c$, so
  $c$ is also an axis for $G$. Thus every element of the subgroup $H$
  generated by $F,G$ has $c$ as an axis, giving a homomorphism $H\to\mathbb{R}$
  sending an isometry to its period along $c$, with trivial kernel (a
  nontrivial deck transformation cannot fix $c$ pointwise). The image
  $A\subset\mathbb{R}$ is an additive subgroup with $a=\inf\{x\in A\mid x>0\}>0$,
  since no nonzero period along $c$ can be smaller than
  $\operatorname{inj}(M)/|\dot c|$ by
  \cref{lem:pet-ch6-deck-transformation-dilation}; a nonzero-infimum additive
  subgroup of $\mathbb{R}$ is cyclic, $A=\{na\mid n\in\mathbb{Z}\}$, so $H$ (and hence
  any Abelian subgroup of $\pi_1(M)$, which is generated by commuting
  elements) is cyclic. A compact product $M\times N$ has $\pi_1(M\times
  N)=\pi_1(M)\times\pi_1(N)$, which is Abelian but not cyclic unless one
  factor is simply connected; if $\dim M,\dim N\ge1$ neither factor is
  contractible in a compact manifold with a metric, so $\pi_1(M)\times
  \pi_1(N)$ fails to be cyclic in general, precluding negative curvature.
  \source{petersen:page-0265}
\end{proof}

\section{Positive Curvature}

\begin{definition}[sectional curvature bounds]
  \label{def:pet-ch6-sec-bounds}
  \uses{def:pet-ch3-sectional-curvature}
  \lean{PetersenLib.HasSecBoundedBelow,PetersenLib.HasSecBoundedAbove,
    PetersenLib.HasSecIn,PetersenLib.HasSecPos,PetersenLib.HasSecNeg}
  \leanok
  Let $(M,g)$ be a Riemannian manifold with Levi-Civita connection $D$ and let
  $k,K\in\mathbb{R}$. Say $(M,g)$ has $\sec\ge k$ if $\sec(v,w)\ge k$ for every
  $p\in M$ and every linearly independent pair $v,w\in T_pM$; $\sec\le K$ if
  $\sec(v,w)\le K$ for every such $p,v,w$; $k\le\sec\le K$ if both hold; and
  $\sec>0$ (resp.\ $\sec<0$) if $\sec(v,w)>0$ (resp.\ $<0$) for every such
  $p,v,w$. As $\sec$ depends only on the $2$-plane $\operatorname{span}\{v,w\}$,
  these are conditions on the $2$-planes of $TM$.

  This is the vocabulary of every comparison theorem of this section and of
  \emph{Basic Comparison Estimates} below, and it is stated here because
  Petersen uses the notation throughout without ever displaying it as a
  definition.

  Two points are forced rather than chosen. The linear independence hypothesis
  is essential: $\sec(v,w)$ is a quotient by $g(v\wedge w,v\wedge w)$, which
  vanishes exactly on dependent pairs, so dropping it would make $\sec\ge k$
  unsatisfiable for every $k>0$ and every downstream theorem vacuous. And
  $\sec>0$ is \emph{not} $\exists k>0,\ \sec\ge k$: on a noncompact manifold
  positive curvature may decay to $0$ at infinity (a paraboloid) with no uniform
  positive lower bound, so \Cref{thm:pet-ch6-synge-theorem} needs its own
  predicate.
\end{definition}

\begin{definition}[Ricci curvature lower bound]
  \label{def:pet-ch6-ricci-lower-bound}
  \lean{PetersenLib.HasRicciBoundedBelowAt,PetersenLib.HasRicciBoundedBelow}
  \uses{def:pet-ch3-ricci-curvature}
  \leanok
  Let $n=\dim M$.  Say that $(M,g)$ has
  $\operatorname{Ric}\ge(n-1)k\,g$ if, for every $p\in M$ and $v\in T_pM$,
  \[
    (n-1)k\,g(v,v)\le \operatorname{Ric}(v,v).
  \]
  In particular, for a unit vector $v$ this reads
  $\operatorname{Ric}(v,v)\ge(n-1)k$.
\end{definition}

\begin{lemma}[Bonnet--Synge index inequality, scalar core]
  \label{lem:pet-ch6-bonnet-synge-index-core}
  \lean{PetersenLib.bonnetSynge_index_core,PetersenLib.integral_sin_sq_window,
    PetersenLib.integral_cos_sq_window,PetersenLib.sq_pi_div_lt_of_pi_div_sqrt_lt}
  \uses{def:pet-ch6-sec-bounds}
  \leanok
  Fix $l>0$, $k>0$ with $l>\pi/\sqrt k$, and let $\kappa:[0,l]\to\mathbb R$ satisfy
  $\kappa\ge k$ on $[0,l]$, with $\sin^2(\pi\cdot/l)\,\kappa$ integrable. Then
  \[
    \Big(\frac\pi l\Big)^2\int_0^l\cos^2\Big(\frac\pi lt\Big)dt
      -\int_0^l\sin^2\Big(\frac\pi lt\Big)\kappa(t)\,dt<0 .
  \]
  This is the pure-analysis heart of Bonnet--Synge (and, summed over an orthonormal frame,
  of Myers, \cref{thm:pet-ch6-myers-ricci-diameter}); the ingredients are the window
  integrals $\int_0^l\sin^2(\pi t/l)\,dt=\int_0^l\cos^2(\pi t/l)\,dt=l/2$ and the equivalence
  of $l>\pi/\sqrt k$ with $(\pi/l)^2<k$.
  \source{petersen:page-0266,petersen:page-0267}
\end{lemma}
\begin{proof}
  \leanok
  Both trigonometric integrals equal $l/2$ (change variables $x=\pi t/l$, which sends
  $[0,l]$ onto $[0,\pi]$ and reduces to $\int_0^\pi\sin^2=\int_0^\pi\cos^2=\pi/2$). Since
  $\kappa\ge k$ and $\sin^2\ge0$, $\int_0^l\sin^2(\pi t/l)\kappa\ge k\int_0^l\sin^2(\pi t/l)
  =kl/2$, so the left side is $\le(l/2)\big((\pi/l)^2-k\big)<0$ because $(\pi/l)^2<k$.
\end{proof}

\begin{lemma}[Lemma 6.3.1 — Bonnet, 1855, and Synge, 1926]
  \label{lem:pet-ch6-bonnet-synge-diameter}
  \lean{PetersenLib.bonnetSynge_longGeodesicsNotMinimizing,
    PetersenLib.bonnetSynge_variation_integrand}
  \uses{rem:pet-ch6-second-variation-special-cases,def:pet-ch6-sec-bounds,%
    lem:pet-ch6-bonnet-synge-index-core}
  \leanok
  If $(M,g)$ satisfies $\sec\ge k>0$, then geodesics of length $>\pi/\sqrt k$
  cannot be locally minimizing.

  \emph{Formalization note.} What is mechanized is Petersen's computation itself: for a smooth
  proper variation $f$ of the unit-speed geodesic $c=f(0,\cdot)$ whose variation field is
  $\sin(\pi\cdot/l)E$ with $E$ a unit parallel field $\perp\dot c$, the second variation
  $\tfrac{d^2}{ds^2}E(f_s)|_{0}$ is \emph{strictly negative}. The pointwise Synge integrand is
  evaluated to $(\pi/l)^2\cos^2(\pi t/l)-\sin^2(\pi t/l)\sec(E,\dot c)$
  (\texttt{bonnetSynge\_variation\_integrand}, via the covariant Leibniz rule
  $\dot V=(\pi/l)\cos(\pi t/l)E$ and the $\sec$ bridge with $|E\wedge\dot c|^2=1$) and
  integrated against \cref{lem:pet-ch6-bonnet-synge-index-core}. Three inputs enter as
  hypotheses, each a separate matter: the \emph{existence} of such a smooth variation — the
  exponential variation $\bar c(s,t)=\exp_{c(t)}(sV(t))$ realizes the field
  (\texttt{variationField\_expVariation}), but its joint slab-smoothness is the analytic gap
  documented in \texttt{Ch06/ExpVariation.lean}; the integrability of $t\mapsto\sec(E,\dot c)$
  (a smoothness consequence); and the reading of ``not locally minimizing'' off the strictly
  negative second variation of a proper variation.
  \source{petersen:page-0266}
\end{lemma}
\begin{proof}
  \uses{rem:pet-ch6-second-variation-special-cases,lem:pet-ch6-bonnet-synge-index-core}
  Let $c:[0,l]\to M$ be unit speed with $l>\pi/\sqrt k$. Let $E$ be a unit
  parallel field perpendicular to $c$ and $V(t)=\sin(\tfrac\pi lt)E(t)$; $V$
  vanishes at $t=0,l$, giving a proper variation. By
  \cref{rem:pet-ch6-second-variation-special-cases},
  \[
    \frac{d^2E}{ds^2}\bigg|_{s=0}
    =\Big(\frac\pi l\Big)^2\int_0^l\cos^2\Big(\frac\pi lt\Big)dt
    -\int_0^l\sin^2\Big(\frac\pi lt\Big)\sec(E,\dot c)\,dt
    \le\Big(\frac\pi l\Big)^2\int_0^l\cos^2\Big(\frac\pi lt\Big)dt
    -k\int_0^l\sin^2\Big(\frac\pi lt\Big)dt.
  \]
  Since $l>\pi/\sqrt k$, $(\pi/l)^2<k$, so the right side is
  $<k\int_0^l\cos^2(\tfrac\pi lt)dt-k\int_0^l\sin^2(\tfrac\pi lt)dt=0$ (the two
  integrals are equal by symmetry). So nearby curves in the variation are
  strictly shorter than $c$.
  \source{petersen:page-0266,petersen:page-0267}
\end{proof}

\begin{corollary}[Corollary 6.3.2 — Hopf and Rinow, 1931, and Myers, 1932]
  \label{cor:pet-ch6-hopf-rinow-myers-diameter}
  \lean{PetersenLib.hopfRinowMyers_diameterBound,
    PetersenLib.longGeodesic_notMinimizing_of_secBound,
    PetersenLib.riemannianDistance_lt_of_negSecondVar_unitSpeed,
    PetersenLib.diameterBound_of_notMinimizing,
    PetersenLib.IsBonnetSyngeVariation}
  \uses{lem:pet-ch6-bonnet-synge-diameter,def:pet-ch6-sec-bounds}
  \leanok
  If $(M,g)$ is complete and satisfies $\sec\ge k>0$, then $M$ is compact and
  $\operatorname{diam}(M,g)\le\pi/\sqrt k=\operatorname{diam}S^n_k$; in particular $\pi_1(M)$ is
  finite.

  \emph{Formalization note (compactness + diameter bound; $\pi_1$-finiteness omitted).} The
  mechanized theorem \texttt{hopfRinowMyers\_diameterBound} takes a \emph{genuine} curvature
  hypothesis $\sec\ge k>0$ (\texttt{HasSecBoundedBelow}) together with the \emph{technical}
  existence hypothesis \texttt{hvar}: for every unit-speed distance-realizing segment $\sigma$
  longer than $\pi/\sqrt k$ there exists a Bonnet--Synge variation
  (\texttt{IsBonnetSyngeVariation}) --- a proper exponential variation with field $\sin(\pi
  t/l)E$, $E$ a unit parallel field $\perp\dot\sigma$. This \texttt{hvar} is \emph{exactly} the
  joint slab-smoothness of the exponential variation already carried as a hypothesis by
  \cref{lem:pet-ch6-bonnet-synge-diameter} (which is \texttt{\textbackslash leanok}), here
  existentially quantified over segments; it asserts \emph{regularity}, not the not-minimizing
  conclusion. The curvature hypothesis enters genuinely:
  \texttt{longGeodesic\_notMinimizing\_of\_secBound} feeds $\sec\ge k$ through
  \texttt{bonnetSynge\_longGeodesicsNotMinimizing} (hence \texttt{bonnetSynge\_index\_core}) to
  obtain $d^2E/ds^2|_0<0$, then the packaged bridge
  \texttt{riemannianDistance\_lt\_of\_negSecondVar\_unitSpeed} --- unit-speed base facts $E=l/2$,
  $L=l$, first-variation differentiability, slice smoothness, plus the index-form engine
  \texttt{riemannianDistance\_lt\_of\_secondVariation\_neg} /
  \texttt{isLocalMin\_deriv\_deriv\_nonneg} (a mathlib gap) --- yields $d(\sigma(0),\sigma(l))<l$,
  contradicting distance-realization. The diameter-and-compactness passage
  (\texttt{diameterBound\_of\_notMinimizing}) is Hopf--Rinow segments
  (\texttt{completeManifold\_exists\_isSegment\_unitSpeed}) plus Heine--Borel (Ch.~5,
  \texttt{hopfRinowTheorem}). The finiteness of $\pi_1(M)$ is \emph{not} formalized: it needs
  covering-space theory absent from the development.
  \source{petersen:page-0267}
\end{corollary}
\begin{proof}
  \uses{lem:pet-ch6-bonnet-synge-diameter}
  By \cref{lem:pet-ch6-bonnet-synge-diameter}, no geodesic of length
  $>\pi/\sqrt k$ realizes the distance between its endpoints; since $M$ is
  complete every pair of points is joined by a minimizing geodesic, so
  $\operatorname{diam}(M,g)\le\pi/\sqrt k$. In particular $M$ is bounded and complete,
  hence compact (Hopf--Rinow). The universal cover of $M$ satisfies the same
  curvature hypothesis and is likewise complete, hence compact; a compact
  cover forces the deck group $\pi_1(M)$ to be finite.
  \source{petersen:page-0267}
\end{proof}

\begin{lemma}[Myers' second-variation core]
  \label{lem:pet-ch6-myers-index}
  \lean{PetersenLib.myersRicci_secondVariation_neg,PetersenLib.bonnetSynge_secondVariation_eq,
    PetersenLib.myers_index_core}
  \uses{lem:pet-ch6-bonnet-synge-index-core,rem:pet-ch6-second-variation-special-cases,def:pet-ch6-sec-bounds}
  \leanok
  Let $c:[0,l]\to M$ be a unit-speed geodesic with $l>\pi/\sqrt k$, $k>0$, and let
  $E_1=\dot c,E_2,\dots,E_n$ be a $g$-orthonormal parallel frame along $c$ with
  $\operatorname{Ric}(\dot c,\dot c)=\sum_{i=2}^n\sec(E_i,\dot c)\ge(n-1)k$. Then, taking the
  $m=n-1$ perpendicular fields $V_i(t)=\sin(\pi t/l)E_i(t)$, at least one has a strictly
  negative second variation $\tfrac{d^2}{ds^2}E(f_{i,s})|_0<0$; hence $c$ is not locally
  minimizing.

  \emph{Formalization note.} This is the analytic heart of \cref{thm:pet-ch6-myers-ricci-diameter}.
  Each perpendicular field gives second variation
  $\tfrac{d^2}{ds^2}E(f_{i,s})|_0=(\pi/l)^2(l/2)-\int_0^l\sin^2(\pi t/l)\sec(E_i,\dot c)\,dt$
  (\texttt{bonnetSynge\_secondVariation\_eq}, the \cref{lem:pet-ch6-bonnet-synge-diameter}
  computation stopped one step early to keep the value); summing over the frame and using the
  frame-sum Ricci bound $(n-1)k\le\sum_i\sec(E_i,\dot c)$ makes the total strictly negative
  (\texttt{myers\_index\_core}), so some summand is $<0$. As in
  \cref{lem:pet-ch6-bonnet-synge-diameter} the variations, their slab-smoothness, the frame's
  parallelism/orthonormality, and the integrability are hypotheses. The Ricci bound is carried in
  its frame-sum form; it equals $\operatorname{Ric}(\dot c,\dot c)\ge(n-1)k$ by the trace identity
  $\operatorname{Ric}(\dot c,\dot c)=\sum_i\sec(E_i,\dot c)$
  (\texttt{ricciCurvature\_eq\_sum\_sectionalCurvature}, §3.1.4) for the orthonormal frame.
  \source{petersen:page-0267,petersen:page-0268}
\end{lemma}
\begin{proof}
  \uses{lem:pet-ch6-bonnet-synge-index-core,rem:pet-ch6-second-variation-special-cases}
  \leanok
  For each $i=2,\dots,n$, $V_i$ vanishes at $t=0,l$ (a proper variation), and Synge's second
  variation formula (\cref{rem:pet-ch6-second-variation-special-cases}) evaluates to
  $(\pi/l)^2(l/2)-\int_0^l\sin^2(\pi t/l)\sec(E_i,\dot c)\,dt$. Summing over the $m=n-1$
  perpendicular directions and pulling the finite sum inside the integral gives
  $(n-1)(\pi/l)^2(l/2)-\int_0^l\sin^2(\pi t/l)\big(\sum_i\sec(E_i,\dot c)\big)dt$. With
  $\sum_i\sec(E_i,\dot c)=\operatorname{Ric}(\dot c,\dot c)\ge(n-1)k$ and $(\pi/l)^2<k$, the
  scalar bound \cref{lem:pet-ch6-bonnet-synge-index-core} (applied $n-1$ times) makes the total
  $<(n-1)(l/2)\big((\pi/l)^2-k\big)<0$. A finite sum of reals that is negative has a negative
  summand, so some $V_i$ gives $\tfrac{d^2}{ds^2}E(f_{i,s})|_0<0$, and $c$ is not locally
  minimizing.
  \source{petersen:page-0267,petersen:page-0268}
\end{proof}

\begin{theorem}[Theorem 6.3.3 — Myers, 1941]
  \label{thm:pet-ch6-myers-ricci-diameter}
  \notready
  \uses{lem:pet-ch6-myers-index,cor:pet-ch6-hopf-rinow-myers-diameter,
    def:pet-ch6-ricci-lower-bound}
  If $(M,g)$ is a complete Riemannian manifold with $\operatorname{Ric}\ge(n-1)k>0$,
  then $\operatorname{diam}(M,g)\le\pi/\sqrt k$, and $(M,g)$ has finite fundamental
  group.
  \source{petersen:page-0267}
\end{theorem}
\begin{proof}
  \uses{rem:pet-ch6-second-variation-special-cases}
  As in \cref{cor:pet-ch6-hopf-rinow-myers-diameter}, it suffices to show no
  geodesic $c:[0,l]\to M$ of length $l>\pi/\sqrt k$ is minimal. Let
  $\dot c(0)=E_1,E_2,\dots,E_n$ be an orthonormal basis of $T_{c(0)}M$, extend
  $E_2,\dots,E_n$ by parallel translation, and set
  $V_i(t)=\sin(\tfrac\pi lt)E_i(t)$, $i=2,\dots,n$. Summing the second
  variation formula (\cref{rem:pet-ch6-second-variation-special-cases}) over
  $i$,
  \begin{align*}
    \sum_{i=2}^n\frac{d^2E}{ds^2}\bigg|_{s=0}
    &=(n-1)\Big(\frac\pi l\Big)^2\int_0^l\cos^2\Big(\frac\pi lt\Big)dt
    -\int_0^l\sin^2\Big(\frac\pi lt\Big)\operatorname{Ric}(\dot c,\dot c)\,dt\\
    &<(n-1)k\int_0^l\cos^2\Big(\frac\pi lt\Big)dt-(n-1)k\int_0^l\sin^2\Big(\frac\pi lt\Big)dt=0,
  \end{align*}
  using $\sum_{i=2}^ng(R(E_i,\dot c)\dot c,E_i)=\operatorname{Ric}(\dot c,\dot c)\ge(n-1)k$
  and $(\pi/l)^2<k$ as before. So at least one $V_i$ gives a negative second
  variation, and $c$ is not locally minimizing. The diameter and finiteness
  of $\pi_1(M)$ then follow exactly as in
  \cref{cor:pet-ch6-hopf-rinow-myers-diameter}.
  \source{petersen:page-0267,petersen:page-0268}
\end{proof}

\begin{example}[Example 6.3.4 — punctured sphere]
  \label{ex:pet-ch6-punctured-sphere-infinite-pi1}
  \lean{PetersenLib.puncturedSphereExample}
  The incomplete manifold $S^2-\{\pm p\}$ has constant curvature $1$ and
  infinite fundamental group; its universal cover also has diameter $\pi$.
  This shows completeness cannot be dropped in
  \cref{cor:pet-ch6-hopf-rinow-myers-diameter,thm:pet-ch6-myers-ricci-diameter}.
  \source{petersen:page-0268}
\end{example}

\begin{example}[Example 6.3.5 — doubly warped product with $\operatorname{Ric}>0$]
  \label{ex:pet-ch6-doubly-warped-ricci-positive}
  \lean{PetersenLib.dwpRicciPositiveOnS1TimesR3}
  \uses{thm:pet-ch6-myers-ricci-diameter}
  $S^1\times\mathbb{R}^3$ admits a complete doubly warped product metric
  $dr^2+\rho^2(r)d\theta^2+\phi^2(r)ds_r^2$ with $\operatorname{Ric}>0$ everywhere,
  showing incompleteness is not needed to violate the topological conclusions
  of \cref{thm:pet-ch6-myers-ricci-diameter} on noncompact manifolds — Myers'
  theorem itself forces compactness once $\operatorname{Ric}$ is bounded below by a
  positive constant, but here $\operatorname{Ric}$ is merely positive (not bounded
  away from $0$). With $\rho(t)=t^{-1/4}$, $\phi(t)=t^{3/4}$ for $t>1$ the
  Ricci curvature is positive; extending to $[0,\infty)$ so the metric is
  $C^\infty$ at $0$, one may arrange $\rho,\phi$ to be $C^1$ and piecewise
  smooth at $t=1$ with $-1<\dot\rho\le0$, $0<\dot\phi\le1$, and, on $[0,1]$,
  $\dot\rho\le0$; smoothing $\rho$ at $t=1$ while preserving positivity of
  $\operatorname{Ric}$ then yields the example.
  \source{petersen:page-0268}
\end{example}

\begin{remark}[closed parallel fields around closed geodesics]
  \label{rem:pet-ch6-closed-parallel-field-construction}
  \lean{PetersenLib.closedParallelFieldAroundClosedGeodesic}
  \uses{def:pet-ch6-parallel-field}
  Let $c:[0,l]\to M$ be a closed unit speed geodesic, $p=c(0)=c(l)$, and
  $P:T_pM\to T_pM$ parallel translation once around $c$, a linear isometry
  fixing $\dot c(0)$ (since $P(\dot c(0))=\dot c(l)=\dot c(0)$), hence
  preserving $\dot c(0)^\perp\subset T_pM$. A linear isometry
  $L:\mathbb{R}^k\to\mathbb{R}^k$ with $\det L=(-1)^{k+1}$ has $1$ as an eigenvalue.
  Consequently: (1) if $M$ is orientable and even-dimensional, $P$ preserves
  orientation ($\det P=1$) and $\dot c(0)^\perp$ is odd-dimensional, so
  $P|_{\dot c(0)^\perp}$ has $\det=(-1)^{(\dim M-1)+1}=1$... more precisely,
  since $\dim(\dot c(0)^\perp)=\dim M-1$ is odd, $P|_{\dot c(0)^\perp}$ fixes
  a vector, giving a closed parallel field around $c$; (2) if $M$ is
  nonorientable, odd-dimensional, and $c$ is orientation-reversing, then
  $\dot c(0)^\perp$ is even-dimensional and $P|_{\dot c(0)^\perp}$ has
  $\det=-1$, again fixing a vector and yielding a closed parallel field.
  \source{petersen:page-0268,petersen:page-0269}
\end{remark}

\begin{theorem}[Theorem 6.3.6 — Synge, 1936]
  \label{thm:pet-ch6-synge-theorem}
  \lean{PetersenLib.syngeTheorem}
  \uses{lem:pet-ch6-deck-transformation-dilation,rem:pet-ch6-closed-parallel-field-construction,rem:pet-ch6-second-variation-special-cases,def:pet-ch6-sec-bounds}
  Let $M$ be a compact manifold with $\sec>0$. (1) If $M$ is even-dimensional
  and orientable, $M$ is simply connected. (2) If $M$ is odd-dimensional,
  $M$ is orientable.
  \source{petersen:page-0269}
\end{theorem}
\begin{proof}
  \uses{lem:pet-ch6-deck-transformation-dilation, rem:pet-ch6-closed-parallel-field-construction, rem:pet-ch6-second-variation-special-cases}
  Argue by contradiction: suppose the universal cover $\pi:\tilde M\to M$ is
  nontrivial, and choose a nontrivial deck transformation $F$, orientation
  reversing along a closed geodesic in the odd-dimensional case. By
  \cref{lem:pet-ch6-deck-transformation-dilation} there is a unit speed axis
  $\tilde c$ for $F$ projecting to $c=\pi\circ\tilde c$, the shortest curve in
  its free homotopy class in $M$, restricted to $[a,b]$ of length
  $b-a=\min\delta_F$; $c$ is a closed geodesic. In both hypothesis cases,
  \cref{rem:pet-ch6-closed-parallel-field-construction} produces a unit
  parallel field $E$ along $c$, closed and perpendicular to $\dot c$. Using
  $E$ as a (non-proper, but with cancelling endpoint terms since $c$ is
  closed) variational field in
  \cref{rem:pet-ch6-second-variation-special-cases},
  \[
    \frac{d^2E(c_s)}{ds^2}\bigg|_{s=0}
    =-\int_a^bg(R(E,\dot c)\dot c,E)\,dt=-\int_a^b\sec(E,\dot c)\,dt<0.
  \]
  So nearby closed curves in this variation are strictly shorter than $c$,
  contradicting minimality of $c$ in its free homotopy class.
  \source{petersen:page-0269}
\end{proof}

\begin{remark}[the Hopf problem]
  \label{rem:pet-ch6-hopf-problem}
  \lean{PetersenLib.hopfProblemRemark}
  \uses{thm:pet-ch6-synge-theorem}
  \Cref{thm:pet-ch6-synge-theorem} shows $\mathbb{RP}^2\times\mathbb{RP}^2$, which has
  positive Ricci curvature, admits no metric of positive sectional curvature
  (it is even-dimensional, orientable, but not simply connected). Whether
  $S^2\times S^2$ admits positive sectional curvature is open (the Hopf
  problem); \cref{thm:pet-ch6-preissmann} ruled out negative curvature on any
  product manifold via fundamental group considerations, but such
  considerations alone do not resolve the positive-curvature case.
  \source{petersen:page-0269}
\end{remark}

\section{Basic Comparison Estimates}

\begin{proposition}[Proposition 6.4.1 — Riccati comparison principle]
  \label{prop:pet-ch6-riccati-comparison-principle}
  \lean{PetersenLib.riccatiComparisonPrinciple}
  \leanok
  Let $\rho_1,\rho_2:(0,b)\to\mathbb{R}$ be smooth with
  \[
    \dot\rho_1+\rho_1^2\le\dot\rho_2+\rho_2^2 ,
  \]
  and let $F$ be an antiderivative of $\rho_1+\rho_2$ on $(0,b)$. Then
  $(\rho_2-\rho_1)e^{F}$ is nondecreasing on $(0,b)$; equivalently, for all
  $0<s\le t<b$,
  \[
    \rho_2(t)-\rho_1(t)\ \ge\ \big(\rho_2(s)-\rho_1(s)\big)
      \exp\Big({-\int_s^t(\rho_1+\rho_2)}\Big).
  \]
  \source{petersen:page-0270}
\end{proposition}
\begin{proof}
  \leanok
  Since $\dot F=\rho_1+\rho_2$, the product rule gives, at every $t\in(0,b)$,
  \[
    \frac{d}{dt}\big((\rho_2-\rho_1)e^{F}\big)
      =(\dot\rho_2-\dot\rho_1)e^{F}+(\rho_2-\rho_1)(\rho_1+\rho_2)e^{F}
      =\big(\dot\rho_2-\dot\rho_1+\rho_2^2-\rho_1^2\big)e^{F},
  \]
  where the last step is $(\rho_2-\rho_1)(\rho_2+\rho_1)=\rho_2^2-\rho_1^2$. The
  hypothesis says exactly that $\dot\rho_2-\dot\rho_1+\rho_2^2-\rho_1^2\ge0$, and
  $e^{F}>0$, so this derivative is $\ge0$ and $(\rho_2-\rho_1)e^{F}$ is
  nondecreasing on the interval $(0,b)$. For $0<s\le t<b$ this gives
  $(\rho_2(t)-\rho_1(t))e^{F(t)}\ge(\rho_2(s)-\rho_1(s))e^{F(s)}$; dividing by
  $e^{F(t)}>0$ yields the displayed inequality, in which the factor
  $e^{F(s)-F(t)}=\exp(-\int_s^t(\rho_1+\rho_2))$ is independent of the choice of
  antiderivative.
  \source{petersen:page-0270}
\end{proof}

\begin{remark}[the printed conclusion of Proposition 6.4.1]
  \label{rem:pet-ch6-riccati-principle-correction}
  \lean{PetersenLib.riccatiComparisonPrinciple_statement_counterexample}
  \leanok
  \uses{prop:pet-ch6-riccati-comparison-principle}
  \emph{Correction.} The printed source states the conclusion of
  Proposition 6.4.1 as
  \[
    \rho_2-\rho_1\ \ge\ \limsup_{t\to0}\big(\rho_2(t)-\rho_1(t)\big),
  \]
  which is false as stated. On $(0,2)$ put
  \[
    \rho_1(t)=\frac{1-e^{-t}}{2},\qquad \rho_2(t)=\frac{1+e^{-t}}{2},
  \]
  so $\rho_1+\rho_2=1$ and $\rho_2-\rho_1=e^{-t}$. Then
  $\rho_2^2-\rho_1^2=(\rho_2-\rho_1)(\rho_2+\rho_1)=e^{-t}$ while
  $\dot\rho_2-\dot\rho_1=-e^{-t}$, so $\dot\rho_1+\rho_1^2=\dot\rho_2+\rho_2^2$
  and the hypothesis holds with equality. Yet
  $\limsup_{t\to0}(\rho_2-\rho_1)=1$ while $\rho_2(1)-\rho_1(1)=e^{-1}<1$; since
  $\rho_2-\rho_1=e^{-t}$ is strictly decreasing, the asserted inequality in fact
  fails at \emph{every} $t\in(0,2)$.

  The gap is in the printed proof, which establishes only that
  $(\rho_2-\rho_1)e^{F}$ is nondecreasing. Deducing the asserted inequality from
  that would require $e^{F(s)-F(t)}\ge1$ in the limit $s\to0^+$, which is not
  automatic: here $F(t)=t$, so $e^{F(s)-F(t)}\to e^{-t}<1$. The monotonicity
  itself is correct, and it is all that
  \cref{cor:pet-ch6-riccati-comparison-estimate} consumes — through
  \cref{lem:pet-ch6-riccati-comparison-le} — so no downstream result is affected.
  \source{petersen:page-0270}
\end{remark}

\begin{lemma}[Riccati comparison — corrected conclusion]
  \label{lem:pet-ch6-riccati-comparison-le}
  \lean{PetersenLib.riccatiComparisonPrinciple_le_of_liminf}
  \leanok
  \uses{prop:pet-ch6-riccati-comparison-principle}
  In the setting of \cref{prop:pet-ch6-riccati-comparison-principle}, suppose in
  addition that
  \[
    \liminf_{t\to0^+}\big(\rho_2(t)-\rho_1(t)\big)e^{F(t)}\ \ge\ 0 .
  \]
  Then $\rho_1\le\rho_2$ on all of $(0,b)$.
\end{lemma}
\begin{proof}
  \leanok
  \uses{prop:pet-ch6-riccati-comparison-principle}
  Write $G=(\rho_2-\rho_1)e^{F}$, which is nondecreasing on $(0,b)$ by
  \cref{prop:pet-ch6-riccati-comparison-principle}. Fix $t\in(0,b)$ and suppose,
  for contradiction, that $G(t)<0$. Every $s\in(0,t)$ satisfies $G(s)\le G(t)$ by
  monotonicity, so $\liminf_{s\to0^+}G(s)\le G(t)<0$, contradicting the
  hypothesis. Hence $G(t)\ge0$, and since $e^{F(t)}>0$ this gives
  $\rho_2(t)-\rho_1(t)\ge0$.
\end{proof}

\begin{remark}[when the printed conclusion is nevertheless valid]
  \label{rem:pet-ch6-riccati-liminf-hypothesis}
  \lean{PetersenLib.riccatiComparisonPrinciple_le_of_bounded_of_exp_tendsto_zero}
  \leanok
  \uses{lem:pet-ch6-riccati-comparison-le}
  If $\rho_2-\rho_1$ is bounded near $0^+$ and $e^{F}\to0$ as $t\to0^+$, then the
  side condition of \cref{lem:pet-ch6-riccati-comparison-le} holds, and hence
  $\rho_1\le\rho_2$ on $(0,b)$. The second hypothesis amounts to
  $\int_{0^+}(\rho_1+\rho_2)=-\infty$. This is the situation of
  \cref{cor:pet-ch6-riccati-comparison-estimate}, where
  $\rho_1(t),\rho_2(t)=\tfrac1t+O(t)$: then $\rho_1+\rho_2=\tfrac2t+O(t)$ gives
  $F(t)=2\log t+O(t^2)$ and $e^{F(t)}=t^2e^{O(t^2)}$, while $\rho_2-\rho_1=O(t)$,
  so $G=(\rho_2-\rho_1)e^{F}=O(t^3)\to0$ and the hypothesis holds. In that regime
  $\limsup_{t\to0}(\rho_2-\rho_1)=0$, so the printed conclusion of
  Proposition 6.4.1 reduces to $\rho_1\le\rho_2$ and is correct.

  Both halves of the hypothesis are needed on each $\rho_i$ separately, not just
  on the sum: for $\rho_{1,2}(t)=\tfrac1t\mp\tfrac1{2t^2}$ one still has
  $\rho_1+\rho_2=\tfrac2t$ and $\dot\rho_1+\rho_1^2=\dot\rho_2+\rho_2^2$, but
  $\rho_2-\rho_1=t^{-2}$ is unbounded near $0$ and
  $\limsup_{t\to0}(\rho_2-\rho_1)=+\infty$.
  \source{petersen:page-0270}
\end{remark}
\begin{proof}
  \leanok
  \uses{lem:pet-ch6-riccati-comparison-le}
  Let $C\ge0$ bound $|\rho_2-\rho_1|$ near $0^+$. Given $\varepsilon>0$, choose
  $s>0$ small enough that both $|\rho_2(s)-\rho_1(s)|\le C$ and
  $e^{F(s)}<\varepsilon/(C+1)$. Then
  \[
    \big(\rho_2(s)-\rho_1(s)\big)e^{F(s)}\ \ge\ -C\,e^{F(s)}
      \ >\ -\frac{C\varepsilon}{C+1}\ >\ -\varepsilon ,
  \]
  using $e^{F(s)}>0$ for the first inequality. As $\varepsilon>0$ was arbitrary
  and such $s$ may be taken arbitrarily close to $0$, this gives
  $\liminf_{t\to0^+}(\rho_2-\rho_1)e^{F}\ge0$, and
  \cref{lem:pet-ch6-riccati-comparison-le} applies.
\end{proof}

\begin{corollary}[Corollary 6.4.2 — Riccati comparison estimate]
  \label{cor:pet-ch6-riccati-comparison-estimate}
  \lean{PetersenLib.riccatiComparisonEstimate}
  \leanok
  \uses{prop:pet-ch6-riccati-comparison-principle,lem:pet-ch6-riccati-comparison-le,
    rem:pet-ch6-riccati-liminf-hypothesis,rem:pet-ch6-riccati-estimate-strictness}
  Let $\rho:(0,b)\to\mathbb{R}$ be smooth with $\rho(t)=\tfrac1t+O(t)$ and $k\in\mathbb{R}$.
  (1) If $\dot\rho+\rho^2\le-k$, then $\rho(t)\le\operatorname{sn}_k'(t)/\operatorname{sn}_k(t)$, and
  moreover $b\le\pi/\sqrt k$ when $k>0$. (2) If $-k\le\dot\rho+\rho^2$, then
  $\operatorname{sn}_k'(t)/\operatorname{sn}_k(t)\le\rho(t)$ for all $t<b$ when $k\le0$, and for
  $t<\min\{b,\pi/\sqrt k\}$ when $k>0$.

  (This matches the printed source. An earlier revision of this node claimed
  Petersen asserts the \emph{strict} bound $b<\pi/\sqrt k$ here; he does not —
  see \cref{rem:pet-ch6-riccati-estimate-strictness}.)
  \source{petersen:page-0270,petersen:page-0271}
\end{corollary}

\begin{remark}[the strictness in Corollary 6.4.2(1): why $\le$ and not $<$]
  \label{rem:pet-ch6-riccati-estimate-strictness}
  \lean{PetersenLib.riccatiComparisonEstimate_lt_isFalse}
  \leanok
  The bound $b\le\pi/\sqrt k$ in \cref{cor:pet-ch6-riccati-comparison-estimate}(1)
  cannot be strengthened to $b<\pi/\sqrt k$. Take $k=1$ and
  $\rho=\operatorname{sn}_1'/\operatorname{sn}_1=\cot$ on $b=\pi=\pi/\sqrt k$. Then $\rho$ is
  smooth on $(0,\pi)$, satisfies $\rho(t)=\tfrac1t+O(t)$, and solves
  $\dot\rho+\rho^2=-1$, so every hypothesis of (1) holds — with equality — yet
  $b=\pi/\sqrt k$. The printed argument rules out $b>\pi/\sqrt k$ (there $\rho$
  would have to blow up before $\pi/\sqrt k$, contradicting smoothness on
  $(0,b)$), which gives exactly $b\le\pi/\sqrt k$; the boundary case is realised,
  as the model $\rho_k$ itself shows.

  \emph{Provenance correction — this is not an error in Petersen.} An earlier
  revision of this remark asserted that ``the printed source concludes
  $b<\pi/\sqrt k$'' and recorded a source error against Petersen. That is
  mistaken. Petersen (p.~255) prints ``Moreover, $b\le\pi/\sqrt k$ when $k>0$'',
  which is correct as stated; verified against
  \texttt{references/petersen/pages/page-0271.png} by vision. The strict
  inequality came from the page \emph{transcription}
  \texttt{references/petersen/tex/page-0271.tex}, line~17, which renders the
  ``$\le$'' as ``$<$''. The lemma above remains worth keeping — it pins the
  boundary case and guards the $\le$ against a future ``obvious''
  strengthening — but the attribution was wrong, and the transcription is the
  defect. Nothing downstream depends on the strictness either way.
  \source{petersen:page-0271}
\end{remark}
\begin{proof}
  \uses{lem:pet-ch6-riccati-comparison-le,rem:pet-ch6-riccati-liminf-hypothesis}
  The function $\rho_k(t)=\operatorname{sn}_k'(t)/\operatorname{sn}_k(t)$ solves $\dot\rho_k+\rho_k^2=-k$
  and satisfies $\rho_k(t)=\tfrac1t+O(t)$ for every $k$; when $k>0$ it is
  defined only on $(0,\pi/\sqrt k)$, with
  $\lim_{t\to\pi/\sqrt k}\rho_k(t)=-\infty$. Applying
  \cref{lem:pet-ch6-riccati-comparison-le} to $\rho,\rho_k$ (in the
  order appropriate to each case) gives the stated one-sided inequalities on the
  common domain of definition: both functions are $\tfrac1t+O(t)$, so by
  \cref{rem:pet-ch6-riccati-liminf-hypothesis} the side condition
  $\liminf_{t\to0^+}(\rho_2-\rho_1)e^{F}\ge0$ is met. In case
  (1), if $b>\pi/\sqrt k$ with $k>0$, the inequality $\rho\le\rho_k$ on
  $(0,\pi/\sqrt k)$ together with $\rho_k\to-\infty$ would force $\rho$ to be
  non-smooth (blow up) before reaching $\pi/\sqrt k$, contradicting
  smoothness of $\rho$ on $(0,b)$; hence $b\le\pi/\sqrt k$ (the boundary case
  $b=\pi/\sqrt k$ does occur, so this cannot be improved to a strict inequality —
  see \cref{rem:pet-ch6-riccati-estimate-strictness}). In case (2), the
  comparison function $\rho_k$ is defined only up to $\pi/\sqrt k$ when
  $k>0$, forcing the stated restriction on the range of $t$.
  \source{petersen:page-0271}
\end{proof}

\begin{theorem}[Theorem 6.4.3 — Rauch comparison]
  \label{thm:pet-ch6-rauch-comparison}
  \lean{PetersenLib.rauchComparisonHessian}
  \uses{cor:pet-ch6-riccati-comparison-estimate,rem:pet-ch6-jacobi-hessian-r,def:pet-ch6-sec-bounds}
  Assume $(M,g)$ satisfies $k\le\sec\le K$, and write $g=dr^2+g_r$ for the
  metric in exponential polar coordinates around $p$ on the star-shaped
  region $T_pM-\{0\}$. Then
  \[
    \frac{\operatorname{sn}_K'(r)}{\operatorname{sn}_K(r)}g_r\le\operatorname{Hess}r\le\frac{\operatorname{sn}_k'(r)}{\operatorname{sn}_k(r)}g_r,
  \]
  and consequently the modified distance functions $f_k,f_K$ satisfy
  $\operatorname{Hess}f_k\le(1-kf_k)g$ and $\operatorname{Hess}f_K\ge(1-Kf_K)g$.

  \emph{Correction (transcription, not a source error).} This node previously
  displayed the two comparison functions with $k$ and $K$ interchanged, i.e.\
  $\operatorname{sn}_k'/\operatorname{sn}_k\le\operatorname{Hess}r\le
  \operatorname{sn}_K'/\operatorname{sn}_K$. That inverts the mathematics: the
  \emph{upper} curvature bound $K$ controls the Hessian from \emph{below} and
  the \emph{lower} bound $k$ from above, which is exactly what this node's own
  proof derives ($\dot\rho\ge-K-\rho^2$ yields $\rho\ge
  \operatorname{sn}_K'/\operatorname{sn}_K$ via
  \Cref{cor:pet-ch6-riccati-comparison-estimate}(2)). The swapped form is also
  refutable outright: on flat $\mathbb{R}^2$ ($\sec\equiv0$, so $k=0\le\sec\le K$
  for every $K>0$) it would assert $\operatorname{Hess}r=\tfrac1rg_r\le
  \sqrt K\cot(\sqrt Kr)g_r$, whose right-hand side is negative as soon as
  $r>\pi/(2\sqrt K)$. Verified against the source page by vision: Petersen
  (p.~255) states the corrected form above. Petersen's text is right here; the
  blueprint was wrong.
  \source{petersen:page-0271}
\end{theorem}
\begin{proof}
  \uses{cor:pet-ch6-riccati-comparison-estimate, rem:pet-ch6-jacobi-hessian-r}
  Fix a unit speed geodesic $c(t)$, $c(0)=p$, along which $r$ is smooth on
  $[0,b]$. \emph{Lower bound on $\operatorname{Hess}r$, from $\sec\le K$.}
  For a Jacobi field $J$ along $c$ with $J(0)=0$,
  $\dot J(0)\perp\dot c(0)$ (so $J\perp\dot c$ throughout), set
  $\rho(t)=g(J,\dot J)/|J|^2=\operatorname{Hess}r(J/|J|,J/|J|)$ (by
  \cref{rem:pet-ch6-jacobi-hessian-r}). Then
  \[
    \dot\rho=\frac{-R(J,\dot c,\dot c,J)|J|^2+|\dot J|^2|J|^2-2(g(J,\dot J))^2}{|J|^4}
    \ge-K+\frac{|\dot J|^2|J|^2-2(g(J,\dot J))^2}{|J|^4}\ge-K-\rho^2,
  \]
  the last step by Cauchy--Schwarz.
  \emph{Upper bound on $\operatorname{Hess}r$, from $\sec\ge k$.} Let $E$ be a unit
  parallel field along $c$ perpendicular to $\dot c$, $\rho=g(S(E),E)=
  \operatorname{Hess}r(E,E)$ for the shape operator $S=\nabla\nabla r$; by the Riccati
  equation for $S$ (part (2) of proposition 3.2.11 of the source),
  $\dot\rho=-R(E,\dot c,\dot c,E)-g(S(E),S(E))\le-k-(g(S(E),E))^2=-k-\rho^2$.
  In both cases $\rho(t)=\tfrac1t+O(t)$ as $t\to0$, so
  \cref{cor:pet-ch6-riccati-comparison-estimate} gives the two Hessian
  bounds on $g_r$; the bounds for the modified distance functions follow
  immediately from $\operatorname{Hess}f_k=(1-kf_k')\,\cdot$-type identities relating
  $f_k$ to $r$ (corollary 4.3.4 of the source).
  \source{petersen:page-0272}
\end{proof}

\begin{remark}[Remark 6.4.4 — the index-form proof technique]
  \label{rem:pet-ch6-index-form-remark}
  \lean{PetersenLib.rauchComparisonViaIndexForm}
  \uses{thm:pet-ch6-rauch-comparison}
  A more traditional proof of the lower-bound half of
  \cref{thm:pet-ch6-rauch-comparison} uses the index form (exercise 6.7.25),
  and can also be adapted to the upper-bound case.
  \source{petersen:page-0272}
\end{remark}

\begin{example}[Example 6.4.5 — conjugate points in $S^n_K$]
  \label{ex:pet-ch6-conjugate-point-space-form}
  \lean{PetersenLib.conjugatePointSpaceFormExample,
    PetersenLib.isLeast_positive_zero_snFunction}
  \leanok
  On $S^n_K$, $K>0$, polar coordinates around $p$ give
  $dr^2+\operatorname{sn}_K^2(r)\,d\sigma_{n-1}^2$; at distance $\pi/\sqrt K$ from $p$,
  every direction hits a conjugate point.

  The formalized content is the analytic core: $\operatorname{sn}_K$ solves the
  Jacobi equation $\ddot x+Kx=0$, and $\pi/\sqrt K$ is its \emph{first} positive
  zero (its full zero set being $\{n\pi/\sqrt K\}$), so the sphere directions of
  the polar form degenerate exactly there and nowhere earlier — which is what
  ``every direction hits a conjugate point at distance $\pi/\sqrt K$'' says. The
  companion fact $\operatorname{sn}_K(t)\ne0$ for $t>0$ when $K\le0$ is the
  analytic content of \cref{rem:pet-ch6-injectivity-radius-nonpositive}
  (no conjugate points in nonpositive curvature).
  \source{petersen:page-0273}
\end{example}

\begin{theorem}[Theorem 6.4.6]
  \label{thm:pet-ch6-conjugate-radius-bound}
  \lean{PetersenLib.conjugateRadiusLowerBound}
  \uses{thm:pet-ch6-rauch-comparison,ex:pet-ch6-conjugate-point-space-form}
  If $(M,g)$ has $\sec\le K$, $K>0$, then
  $\exp_p:B(0,\pi/\sqrt K)\to M$ has no critical points.
  \source{petersen:page-0273}
\end{theorem}
\begin{proof}
  \uses{thm:pet-ch6-rauch-comparison, rem:pet-ch6-jacobi-hessian-r}
  Let $c(t)=\exp_p(tv)$ be unit speed and $J$ a Jacobi field along $c$ with
  $J(0)=0$, $\dot J(0)\perp\dot c(0)$; suppose toward a contradiction that
  $J>0$ on $(0,b)$ with $J(b)=0$ for some $b\le\pi/\sqrt K$. From the proof of
  \cref{thm:pet-ch6-rauch-comparison} (upper bound case with $k=K$),
  $g(\dot J(t),J(t))/|J(t)|^2\ge\operatorname{sn}_K'(t)/\operatorname{sn}_K(t)$ for
  $t<\min\{b,\pi/\sqrt K\}$, equivalently
  $\frac{d}{dt}\big(|J(t)|/\operatorname{sn}_K(t)\big)\ge0$. Since $D\exp_p$ is the
  identity at the origin, $|J(t)|=t|\dot J(0)|+O(t^2)$, so by l'H\^opital,
  \[
    \lim_{t\to0}\frac{|J(t)|}{\operatorname{sn}_K(t)}
    =\lim_{t\to0}\frac{g(\dot J(t),J(t))}{|J(t)|}
    =\lim_{t\to0}\frac{\operatorname{Hess}r(J(t),J(t))}{|J(t)|}
    =\lim_{t\to0}\frac{t^{-1}|J(t)|^2+O(t^3)}{|J(t)|}=|\dot J(0)|.
  \]
  Hence $|J(t)|\ge|\dot J(0)|\operatorname{sn}_K(t)>0$ for
  $t<\min\{b,\pi/\sqrt K\}\le\pi/\sqrt K$, contradicting $J(b)=0$ with
  $b\le\pi/\sqrt K$ (taking $b<\pi/\sqrt K$ forces the contradiction directly;
  $b=\pi/\sqrt K$ is excluded since $J(t)>0$ persists up to that limit).
  \source{petersen:page-0273}
\end{proof}

\begin{remark}[injectivity radius in nonpositive curvature]
  \label{rem:pet-ch6-injectivity-radius-nonpositive}
  \lean{PetersenLib.injectivityRadius_nonpositiveCurvature}
  \uses{thm:pet-ch6-cartan-hadamard}
  If $\sec\le0$ there are no conjugate points, so
  $\operatorname{inj}(p)=\tfrac12\cdot(\text{length of shortest geodesic loop at }p)$.
  On a closed manifold with $\sec\le0$,
  $\operatorname{inj}(M)=\inf_p\operatorname{inj}(p)=\tfrac12\cdot(\text{length of shortest closed
  geodesic})$: the infimum is attained by compactness and continuity of
  $p\mapsto\operatorname{inj}(p)$; if $p$ realizes it via a geodesic loop $c:[0,1]\to M$,
  splitting $c$ at its midpoint $c(\tfrac12)$ gives
  $\operatorname{inj}(c(\tfrac12))\le\operatorname{inj}(p)$, forcing $c$ to also be a geodesic loop
  from $c(\tfrac12)$, hence smooth at $p$ and a closed geodesic.
  \source{petersen:page-0274}
\end{remark}

\begin{lemma}[Lemma 6.4.7 — Klingenberg]
  \label{lem:pet-ch6-klingenberg-injectivity-estimate}
  \lean{PetersenLib.klingenbergInjectivityEstimate}
  \uses{thm:pet-ch6-conjugate-radius-bound,rem:pet-ch6-injectivity-radius-nonpositive}
  Let $(M,g)$ be compact with $\sec\le K$, $K>0$. Then
  \[
    \operatorname{inj}(p)\ge\min\Big\{\frac{\pi}{\sqrt K},\ \tfrac12\cdot
    (\text{length of shortest geodesic loop at }p)\Big\},
  \]
  and $\operatorname{inj}(M)\ge\pi/\sqrt K$, or
  $\operatorname{inj}(M)=\tfrac12\cdot(\text{length of shortest closed geodesic})$.
  \source{petersen:page-0274}
\end{lemma}

\begin{theorem}[Theorem 6.4.8]
  \label{thm:pet-ch6-convexity-radius-criterion}
  \lean{PetersenLib.convexityRadiusCriterion}
  \uses{thm:pet-ch6-rauch-comparison}
  Suppose $R>0$ satisfies (1) $R\le\tfrac12\operatorname{inj}(x)$ for all $x\in B(p,R)$,
  and (2) $R\le\tfrac12\cdot\pi/\sqrt K$, where
  $K=\sup\{\sec(\pi)\mid\pi\subset T_xM,\,x\in B(p,R)\}$. Then $r(x)=|xp|$ is
  convex on $B(p,R)$, and any two points of $B(p,R)$ are joined by a unique
  segment lying in $B(p,R)$.
  \source{petersen:page-0274}
\end{theorem}
\begin{proof}
  \uses{thm:pet-ch6-rauch-comparison}
  Condition (1) ensures any two points of $B(p,R)$ are joined by a unique
  segment in $M$, and that $r$ is smooth on $B(p,2R)-\{p\}$; condition (2)
  and \cref{thm:pet-ch6-rauch-comparison} give $\operatorname{Hess}r\ge0$ on $B(p,R)$.
  For $x\in B(p,R)$ fixed, let $C_x=\{y\in B(p,R):\text{the segment
  }[x,y]\subset B(p,R)\}$; $x\in C_x$ and $C_x$ is open by continuity of the
  segment on its endpoints. If $y\in B(p,R)\cap\partial C_x$, the segment
  $c:[0,1]\to M$ from $x$ to $y$ lies in $\overline{B(p,R)}$ by continuity;
  $\varphi(t)=r(c(t))$ satisfies $\varphi(0),\varphi(1)<R$ and
  $\ddot\varphi=\operatorname{Hess}r(\dot c,\dot c)\ge0$, so $\varphi$ is convex and
  $\max\varphi\le\max\{\varphi(0),\varphi(1)\}<R$, forcing $c\subset B(p,R)$,
  i.e.\ $y\in C_x$; thus $C_x$ is also closed in $B(p,R)$, hence (by
  connectedness of $B(p,R)$) $C_x=B(p,R)$.
  \source{petersen:page-0274,petersen:page-0275}
\end{proof}

\begin{definition}[convexity radius]
  \label{def:pet-ch6-convexity-radius}
  \lean{PetersenLib.convexityRadius}
  \leanok
  \uses{def:pet-ch6-convex-function,def:pet-ch5-segment,def:pet-ch5-metric-balls}
  The \emph{convexity radius} at $p$, $\operatorname{conv.rad}(p)$, is the largest
  $R>0$ such that $r(x)=|xp|$ is convex on $B(p,R)$ and any two points in $B(p,R)$
  are joined by a unique segment lying in $B(p,R)$. Globally
  $\operatorname{conv.rad}(M,g)=\inf_p\operatorname{conv.rad}(p)$.

  \emph{Formalization notes.} This is Petersen's own definition --- the two
  geometric properties themselves, not the sufficient hypotheses of
  \cref{thm:pet-ch6-convexity-radius-criterion} ($R\le\tfrac12\operatorname{inj}$,
  $R\le\tfrac12\pi/\sqrt K$). Encoding those hypotheses would define a smaller
  quantity, misname it, and needlessly carry a connection and $\sqrt K$
  junk-value handling; the $\ge\min\{\operatorname{inj}/2,\pi/(2\sqrt K)\}$
  estimate (the hypotheses-imply-conclusions direction) is the separate, still
  open \cref{cor:pet-ch6-convexity-radius-bound}. Both defining properties are
  metric, so \texttt{convexityRadius} takes \emph{no} connection argument and has
  the same shape as $\operatorname{inj}$. Two Lean-specific points: the value is an
  $\mathbb{R}_{\ge0}^\infty$ built as a $\sup$ over $\mathrm{ofReal}$-images
  (mirroring $\operatorname{inj}$), and ``unique segment'' is rendered as existence
  plus $\mathrm{EqOn}$ on the parameter interval $[0,|xy|]$, never as $\exists!$
  over $\mathbb{R}\to M$ --- the latter is false for every pair (a segment is
  unconstrained off its interval) and would silently make
  $\operatorname{conv.rad}\equiv0$. See \texttt{UniqueSegmentsInBall},
  \texttt{IsConvexityRadiusWitness}, \texttt{convexityRadius},
  \texttt{globalConvexityRadius} in \texttt{Ch06/ConvexityRadius.lean}.
  \source{petersen:page-0275}
\end{definition}

\begin{corollary}[convexity-radius lower bound]
  \label{cor:pet-ch6-convexity-radius-bound}
  \notready
  \uses{def:pet-ch6-convexity-radius,thm:pet-ch6-convexity-radius-criterion,lem:pet-ch6-klingenberg-injectivity-estimate}
  By \cref{thm:pet-ch6-convexity-radius-criterion,lem:pet-ch6-klingenberg-injectivity-estimate},
  $\operatorname{conv.rad}(M,g)\ge\min\{\operatorname{inj}(M,g)/2,\ \pi/(2\sqrt K)\}$ where
  $K=\sup\sec(M,g)$; in nonpositive curvature this simplifies to
  $\operatorname{conv.rad}(M,g)=\operatorname{inj}(M,g)/2$. (Split out from
  \cref{def:pet-ch6-convexity-radius} so that definition's Lean does not depend on
  the still-open \cref{thm:pet-ch6-convexity-radius-criterion,lem:pet-ch6-klingenberg-injectivity-estimate}.)
  \source{petersen:page-0275}
\end{corollary}

\section{More on Positive Curvature}

\begin{theorem}[Theorem 6.5.1 — Klingenberg, 1959]
  \label{thm:pet-ch6-klingenberg-even-dim-injectivity}
  \lean{PetersenLib.klingenbergEvenDimInjectivity}
  \uses{lem:pet-ch6-klingenberg-injectivity-estimate,thm:pet-ch6-synge-theorem}
  If $(M,g)$ is a compact orientable even-dimensional manifold with
  $0<\sec\le1$, then $\operatorname{inj}(M,g)\ge\pi$. If $M$ is not orientable,
  $\operatorname{inj}(M,g)\ge\pi/2$.
  \source{petersen:page-0276}
\end{theorem}
\begin{proof}
  \uses{lem:pet-ch6-klingenberg-injectivity-estimate, rem:pet-ch6-closed-parallel-field-construction}
  The nonorientable case follows from the orientable case applied to the
  orientation double cover, which then has $\operatorname{inj}\ge\pi$, so the base has
  $\operatorname{inj}\ge\pi/2$. For the orientable case, by
  \cref{lem:pet-ch6-klingenberg-injectivity-estimate} and $K\le1$, if
  $\operatorname{inj}M<\pi$ then $\operatorname{inj}M$ is realized by a closed geodesic
  $c:[0,2\operatorname{inj}M]\to M$, $2\operatorname{inj}M<2\pi$. Since $M$ is orientable and
  even-dimensional, \cref{rem:pet-ch6-closed-parallel-field-construction}
  and the argument in the proof of \cref{thm:pet-ch6-synge-theorem} produce
  for all small $\varepsilon>0$ curves $c_\varepsilon:[0,2\operatorname{inj}M]\to M$
  converging to $c$ as $\varepsilon\to0$ with $L(c_\varepsilon)<L(c)$. As
  $c_\varepsilon\subset B(c_\varepsilon(0),\operatorname{inj}M)$, there is a unique
  segment from $c_\varepsilon(0)$ to any $c_\varepsilon(t)$; let
  $c_\varepsilon(t_\varepsilon)$ be the farthest point on $c_\varepsilon$ from
  $c_\varepsilon(0)$ and $\sigma_\varepsilon$ the (perpendicular at
  $c_\varepsilon(t_\varepsilon)$) segment joining them. As $\varepsilon\to0$,
  $t_\varepsilon\to\operatorname{inj}M$ and $\sigma_\varepsilon$ subconverges to a
  segment from $c(0)$ to $c(\operatorname{inj}M)$ perpendicular to $c$ there. Since the
  conjugate radius is $\ge\pi>\operatorname{inj}M$ and $c$ is a geodesic loop realizing
  the injectivity radius at $c(0)$, there are only two segments from $c(0)$
  to $c(\operatorname{inj}M)$ (namely the two halves of $c$, neither perpendicular to
  $c$ at the interior point $c(\operatorname{inj}M)$ unless $\operatorname{inj}M=\pi$), a
  contradiction with the existence of the perpendicular limiting segment.
  \source{petersen:page-0276}
\end{proof}

\begin{remark}[odd-dimensional counterexamples]
  \label{rem:pet-ch6-klingenberg-odd-dim-counterexample}
  \lean{PetersenLib.klingenbergOddDimCounterexample}
  \uses{thm:pet-ch6-klingenberg-even-dim-injectivity}
  No analogue of \cref{thm:pet-ch6-klingenberg-even-dim-injectivity} holds in
  odd dimensions: the orientable quotients $S^3/\mathbb{Z}_k$ have a closed
  geodesic (image of the Hopf fiber) of length $2\pi/k\to0$; the Berger
  spheres $(S^3,g_\varepsilon)$ have Hopf fiber a closed geodesic of length
  $2\pi\varepsilon$ with curvatures in $[\varepsilon^2,4-3\varepsilon^2]$, so
  after rescaling the upper bound to $1$ the curvatures lie in
  $[\varepsilon^2/(4-3\varepsilon^2),1]$ and the Hopf fiber has length
  $2\pi\varepsilon\sqrt{4-3\varepsilon^2}<2\pi$ once $\varepsilon<1/\sqrt3$,
  with lower curvature bound then below $1/3$.
  \source{petersen:page-0276,petersen:page-0277}
\end{remark}

\begin{definition}[$l$-connected pair]
  \label{def:pet-ch6-l-connected}
  \lean{PetersenLib.IsLConnected}
  $A\subset X$ is \emph{$l$-connected} if the relative homotopy groups
  $\pi_k(X,A)$ vanish for $k<l$; by Hurewicz's theorem the relative homology
  groups $H_k(X,A)$ then also vanish for $k\le l$, so from the long exact
  sequences of the pair, $\pi_k(A)\to\pi_k(X)$ and $H_k(A)\to H_k(X)$ are
  isomorphisms for $k<l$ and surjective for $k=l$.
  \source{petersen:page-0277}
\end{definition}

\begin{definition}[index of a critical point]
  \label{def:pet-ch6-critical-point-index}
  \lean{PetersenLib.criticalPointIndex}
  A critical point $p\in M$ of a smooth function $f:M\to\mathbb{R}$ has \emph{index}
  $m$ if $\operatorname{Hess}f$ is negative definite on an $m$-dimensional subspace of
  $T_pM$. If $m\ge1$, $p$ cannot be a local minimum of $f$, since $f$
  decreases along directions where the Hessian is negative definite.
  \source{petersen:page-0277}
\end{definition}

\begin{theorem}[Theorem 6.5.2]
  \label{thm:pet-ch6-morse-connectivity}
  \lean{PetersenLib.morseTheoreticConnectivity}
  \uses{def:pet-ch6-l-connected,def:pet-ch6-critical-point-index}
  Let $f:M\to\mathbb{R}$ be smooth and proper. If $b$ is not a critical value and
  every critical point in $f^{-1}([a,b])$ has index $\ge m$, then
  $f^{-1}((-\infty,a))\subset f^{-1}((-\infty,b))$ is $(m-1)$-connected.
  \source{petersen:page-0277}
\end{theorem}
\begin{proof}
  If $f^{-1}([a,b])$ has no critical points, the negative gradient flow
  deforms $f^{-1}((-\infty,b))$ onto $f^{-1}((-\infty,a))$. Otherwise cover
  the (compact, by properness) critical set in $f^{-1}([a,b])$ by finitely
  many coordinate boxes $\overline U_i\subset V_i\approx[-1,1]^n$ in which the
  first $m$ coordinates $x^1,\dots,x^m$ span a subspace where $\operatorname{Hess}f$ is
  negative definite at the critical point. Given
  $\phi:N^{k-1}\to f^{-1}([a,b])$ with $k\le m$ (fixed on
  $\partial f^{-1}((-\infty,a))$ if nonempty): (i) on $M-\bigcup U_i$, flow
  along $-\lambda(x)\nabla f|_x$ for $\lambda\ge0$ vanishing exactly on
  $\bigcup U_i$, deforming $\phi$ rel $\overline U_i$ so that
  $\max_{V_i}f\circ\phi$ decreases and is not attained on $\partial V_i$;
  (ii) on each $V_i$, letting $S_i=\{x^1=\cdots=x^m=0\}\subset V_i$, since
  $k\le m$ a transversality argument produces a homotopy $\phi_t$ avoiding
  $S_i$ for $t>0$ and constant outside $V_i$, still with $\max_{V_i}
  f\circ\phi_t$ not attained on $\partial V_i$ for small $t$; (iii) once
  $\phi$ avoids $S_i$, the flow of the radial field
  $\sum_{j=1}^mx^j\partial_j$ on $V_i$ decreases $\max_{V_i}f\circ\phi$ and
  pushes $\phi$ out of $\overline U_i$. Iterating these three deformations
  moves $\phi$ into $f^{-1}((-\infty,a])$, proving $(m-1)$-connectivity.
  \source{petersen:page-0277,petersen:page-0278}
\end{proof}

\begin{definition}[the path space $\Omega_{A,B}(M)$ and index of a geodesic]
  \label{def:pet-ch6-path-space-AB}
  \lean{PetersenLib.pathSpaceAB,PetersenLib.geodesicIndexAB}
  \leanok
  \uses{def:pet-ch6-critical-point-index}
  For $A,B\subset M$, $\Omega_{A,B}(M)=\{c:[0,1]\to M\mid c(0)\in A,\,c(1)\in
  B\}$. When $A,B$ are compact submanifolds, the energy functional
  $E:\Omega_{A,B}(M)\to[0,\infty)$ behaves as a proper smooth function, its
  critical points are geodesics meeting $A,B$ orthogonally at the endpoints,
  and its variational fields along such a geodesic are those tangent to $A,B$
  at the respective endpoints. The \emph{index} of such a geodesic is $\ge k$
  if there is a $k$-dimensional space of such variational fields on which the
  second variation of energy is negative (in the sense of
  \cref{def:pet-ch6-critical-point-index} applied to $E$ on
  $\Omega_{A,B}(M)$).
  \source{petersen:page-0278}
\end{definition}

\begin{theorem}[Theorem 6.5.3]
  \label{thm:pet-ch6-submanifold-connectivity}
  \lean{PetersenLib.submanifoldIndexConnectivity}
  \uses{def:pet-ch6-path-space-AB,thm:pet-ch6-morse-connectivity}
  Let $M$ be a complete Riemannian manifold and $A\subset M$ a compact
  submanifold. If every geodesic in $\Omega_{A,A}(M)$ perpendicular to $A$ at
  both endpoints has index $\ge k$, then $A\subset M$ is $k$-connected.
  \source{petersen:page-0278}
\end{theorem}
\begin{proof}
  \uses{thm:pet-ch6-morse-connectivity}
  Identifying $A=E^{-1}(0)$ inside $\Omega_{A,A}(M)$, the deformation scheme
  of \cref{thm:pet-ch6-morse-connectivity} (carried out in a suitable
  infinite-dimensional or finite-dimensional-approximation setting for
  $\Omega_{A,A}(M)$) shows $A\subset\Omega_{A,A}(M)$ is $(k-1)$-connected.
  Since $\pi_l(\Omega_{A,A}(M),A)=\pi_{l+1}(M,A)$, this gives
  $\pi_{l+1}(M,A)=0$ for $l<k$, i.e.\ $A\subset M$ is $k$-connected.
  \source{petersen:page-0278}
\end{proof}

\begin{theorem}[Theorem 6.5.4 — Berger, 1958]
  \label{thm:pet-ch6-berger-sphere-theorem}
  \lean{PetersenLib.bergerSphereTheorem}
  \uses{thm:pet-ch6-submanifold-connectivity,lem:pet-ch6-bonnet-synge-diameter,def:pet-ch6-l-connected}
  Let $M$ be a closed $n$-manifold with $\sec\ge1$. If $\operatorname{inj}_p>\pi/2$ for
  some $p\in M$, then $M$ is $(n-1)$-connected, hence a homotopy sphere.
  \source{petersen:page-0279}
\end{theorem}
\begin{proof}
  \uses{thm:pet-ch6-submanifold-connectivity, lem:pet-ch6-bonnet-synge-diameter}
  Apply \cref{thm:pet-ch6-submanifold-connectivity} with $A=\{p\}$. Every
  geodesic loop at $p$ is either constant or has length $>\pi$ (since
  $\operatorname{inj}_p>\pi/2$); by \cref{lem:pet-ch6-bonnet-synge-diameter} such a
  geodesic has, for each unit parallel field $E$ perpendicular to it, a
  proper variation with negative second derivative (modifying $E$ by the
  factor $\sin(\pi t/l)$ as in that proof). Since there is an
  $(n-1)$-dimensional space of such parallel fields, the index of the
  geodesic is $\ge n-1$, so $\{p\}\subset M$ is $(n-1)$-connected by
  \cref{thm:pet-ch6-submanifold-connectivity}, i.e.\ $M$ is
  $(n-1)$-connected. Choosing $F:M\to S^n$ of degree $1$, $(n-1)$-connectivity
  of $M$ (and of $S^n$) makes $F$ an isomorphism on $\pi_k$ for $k<n$; by
  Hurewicz, $\pi_n(M)\cong H_n(M)$ and $\pi_n(S^n)\cong H_n(S^n)$, and since
  $\deg F=1$, $F_*:H_n(M)\to H_n(S^n)$ is an isomorphism, so $F_*$ is an
  isomorphism on $\pi_n$ as well. A theorem of Whitehead then shows $F$ is a
  homotopy equivalence.
  \source{petersen:page-0279}
\end{proof}

\begin{theorem}[Theorem 6.5.5 — Klingenberg, 1961]
  \label{thm:pet-ch6-klingenberg-pinching-injectivity}
  \lean{PetersenLib.klingenbergPinchingInjectivity}
  \uses{thm:pet-ch6-klingenberg-even-dim-injectivity,lem:pet-ch6-klingenberg-injectivity-estimate}
  A compact simply connected Riemannian $n$-manifold $(M,g)$ with
  $1\le\sec<4$ has $\operatorname{inj}>\pi/2$.
  \source{petersen:page-0279}
\end{theorem}
\begin{proof}
  \uses{thm:pet-ch6-klingenberg-even-dim-injectivity, lem:pet-ch6-klingenberg-injectivity-estimate}
  By rescaling it suffices to show simply connected manifolds with
  $1<\sec\le4$ have $\operatorname{inj}\ge\pi/2$; the even-dimensional case is
  \cref{thm:pet-ch6-klingenberg-even-dim-injectivity}, so assume $n\ge3$ odd.
  The strict lower curvature bound gives $\delta>0$ such that geodesics of
  length $\ge\pi-\delta$ have index $\ge n-1\ge2$ (as in
  \cref{thm:pet-ch6-berger-sphere-theorem}); hence any loop of constant speed
  based at $p$, viewed in $\Omega_{p,p}(M)$, is homotopic (uniformly, by
  replacing loops with piecewise-geodesic shortenings with breakpoints chosen
  independently of the homotopy parameter) to one of length $<\pi$.

  Suppose toward a contradiction $\operatorname{inj}_p<\pi/2$. By
  \cref{lem:pet-ch6-klingenberg-injectivity-estimate} there is a geodesic
  loop $c_1$ at $p$ of length $<\pi$ realizing $\operatorname{inj}_p$. Simple
  connectivity gives a homotopy of constant-speed loops
  $c_s(t)$, $s\in[0,1]$, $c_s(0)=c_s(1)=p$, all of length $<\pi$, with
  $c_0\equiv p$ and $c_1$ the closed geodesic. Each $c_s\subset B(p,\pi/2)$.
  Pull back the metric on $B(0,\pi/2)\subset T_pM$ so that
  $\exp_p:B(0,\pi/2)\to B(p,\pi/2)$ is a local isometry; each $c_s$ lifts
  uniquely to $\tilde c_s:[0,b_s]\to \tilde B(0,\pi/2)$ with $\tilde c_s(0)=0$,
  where either $\tilde c_s(b_s)\in\partial B(0,\pi/2)$ or $b_s=1$
  (piecewise-geodesic $c_s$ lift by lifting velocities at break points). Let
  $A=\{s: c_s\text{ lifts to a loop }\tilde c_s:[0,1]\to B(0,\pi/2)\text{ at
  }0\}$. Then $0\in A$ (loops near the trivial loop lift near $0$, where
  $\exp_p$ is a diffeomorphism); $A$ is closed, since if $s_i\to s$ in $A$,
  the limit $\tilde c_{s_i}(1)=0$ persists in the limit lift $\tilde c_s$
  (which stays in the interior since $L(c_s)<\pi$); $A$ is open, since for
  $s_0\in A$, choosing $\epsilon>0$ with $\exp_p:B(\tilde
  c_{s_0}(t),\epsilon)\to B(c_{s_0}(t),\epsilon)$ an isometry for all $t$ and
  contained in $B(0,\pi/2)$, nearby loops $c_s\subset\bigcup_tB(c_{s_0}(t),
  \epsilon)$ lift uniquely into $\bigcup_tB(\tilde c_{s_0}(t),\epsilon)$ and
  close up. So $A=[0,1]$; but $c_1$, the closed geodesic of length $<\pi$,
  lifts to a straight line from $0$ that is not a loop (as $\exp_p$
  restricted to the segment is injective on $B(0,\pi/2)$, and a nontrivial
  closed geodesic loop of length $<\pi$ realizing $\operatorname{inj}_p$ cannot close up
  strictly inside $B(0,\pi/2)$ under the local isometry) — a contradiction.
  \source{petersen:page-0279,petersen:page-0280}
\end{proof}

\begin{corollary}[Corollary 6.5.6 — Rauch, Berger, and Klingenberg, 1951--61]
  \label{cor:pet-ch6-rauch-berger-klingenberg-sphere}
  \lean{PetersenLib.rauchBergerKlingenbergSphereTheorem}
  \uses{thm:pet-ch6-berger-sphere-theorem,thm:pet-ch6-klingenberg-pinching-injectivity}
  Let $M$ be a closed simply connected $n$-manifold with $4>\sec\ge1$. Then
  $M$ is $(n-1)$-connected, hence a homotopy sphere.
  \source{petersen:page-0280}
\end{corollary}
\begin{proof}
  \uses{thm:pet-ch6-berger-sphere-theorem, thm:pet-ch6-klingenberg-pinching-injectivity}
  By \cref{thm:pet-ch6-klingenberg-pinching-injectivity}, $\operatorname{inj}(M,g)>\pi/2$
  everywhere; apply \cref{thm:pet-ch6-berger-sphere-theorem} at any $p\in M$.
  \source{petersen:page-0280}
\end{proof}

\begin{corollary}[Corollary 6.5.7]
  \label{cor:pet-ch6-ricci-simply-connected}
  \lean{PetersenLib.ricciDiameterInjectivitySimplyConnected}
  \uses{thm:pet-ch6-submanifold-connectivity,thm:pet-ch6-myers-ricci-diameter}
  If $M$ is a closed $n$-manifold with $\operatorname{Ric}\ge(n-1)$ and $\operatorname{inj}_p>\pi/2$
  for some $p\in M$, then $M$ is simply connected.
  \source{petersen:page-0280}
\end{corollary}
\begin{proof}
  \uses{thm:pet-ch6-submanifold-connectivity, thm:pet-ch6-myers-ricci-diameter}
  As in \cref{thm:pet-ch6-berger-sphere-theorem}, but using the Ricci
  curvature computation of \cref{thm:pet-ch6-myers-ricci-diameter} (summing
  the second variation over an orthonormal frame of parallel fields) instead
  of the sectional curvature bound: every nonconstant geodesic loop at $p$
  has length $>\pi$ and, since $\operatorname{Ric}\ge n-1$, index $\ge1$ (at least one
  parallel-field variation is negative), so by
  \cref{thm:pet-ch6-submanifold-connectivity} with $A=\{p\}$, $\{p\}\subset M$
  is $1$-connected, i.e.\ $M$ is simply connected.
  \source{petersen:page-0280}
\end{proof}

\begin{lemma}[Lemma 6.5.8 — the Connectedness Principle, Wilking, 2003]
  \label{lem:pet-ch6-wilking-connectedness-principle}
  \lean{PetersenLib.wilkingConnectednessPrinciple}
  \uses{def:pet-ch6-l-connected,rem:pet-ch6-second-variation-special-cases}
  Let $M^n$ be compact with $\sec>0$. (a) If $N^{n-k}\subset M^n$ is a closed
  totally geodesic codimension-$k$ submanifold, then $N\subset M$ is
  $(n-2k+1)$-connected. (b) If $N_1^{n-k_1},N_2^{n-k_2}\subset M$ are closed
  totally geodesic submanifolds with $k_1\le k_2$, $k_1+k_2\le n$, then
  $N_1\cap N_2$ is a nonempty totally geodesic submanifold and
  $N_1\cap N_2\to N_2$ is $(n-k_1-k_2)$-connected.
  \source{petersen:page-0280,petersen:page-0281}
\end{lemma}
\begin{proof}
  \uses{thm:pet-ch6-submanifold-connectivity, rem:pet-ch6-second-variation-special-cases}
  \emph{(a)} Let $c\in\Omega_{N,N}(M)$ be a geodesic and $E$ a parallel field
  along $c$ tangent to $N$ at both endpoints. Extend to a variation
  $\tilde c(s,t)$ with $\tilde c(0,t)=c(t)$ and $s\mapsto\tilde c(s,t)$
  geodesics with initial velocity $E|_{c(0)}$; total geodesy of $N$ keeps
  $\tilde c(s,0),\tilde c(s,1)\in N$, so this variation lies in
  $\Omega_{N,N}(M)$. By \cref{rem:pet-ch6-second-variation-special-cases}
  (parallel field, geodesic variation, so both boundary terms vanish),
  $\frac{d^2E(c_s)}{ds^2}|_{s=0}=-\int_0^1g(R(E,\dot c)\dot c,E)\,dt<0$ since
  $\sec>0$ and $E\perp\dot c$. Let $V\subset T_{c(1)}M$ be the set of values
  $E(1)$ over parallel fields $E$ along $c$ with $E(0)\in T_{c(0)}N$; then
  $\dim V=\dim T_{c(0)}N=n-k$. The parallel fields giving negative variations
  correspond to $V\cap T_{c(1)}N$, and since $V,T_{c(1)}N\subset
  \dot c(1)^\perp$ (dimension $n-1$),
  \[
    2n-2k=\dim T_{c(1)}N+\dim V=\dim(V\cap T_{c(1)}N)+\dim(V+T_{c(1)}N)
    \le\dim(V\cap T_{c(1)}N)+n-1,
  \]
  so $\dim(V\cap T_{c(1)}N)\ge n-2k+1$: every geodesic in $\Omega_{N,N}(M)$
  perpendicular to $N$ at both endpoints has index $\ge n-2k+1$, and
  \cref{thm:pet-ch6-submanifold-connectivity} gives $(n-2k+1)$-connectivity
  of $N\subset M$.

  \emph{(b)} $T_p(N_1\cap N_2)=T_pN_1\cap T_pN_2$ for $p\in N_1\cap N_2$ shows
  $N_1\cap N_2$ is totally geodesic. For nonemptiness, take a shortest
  geodesic $c$ from $N_1$ to $N_2$ (perpendicular to both at its endpoints);
  the dimension bound $k_1+k_2\le n$ produces (as in part (a), applied with
  the two submanifolds at the two endpoints) an
  $(n-k_1-k_2+1)$-dimensional space of parallel fields along $c$ tangent to
  $N_1,N_2$ at the endpoints, giving a negative second variation and hence a
  strictly shorter nearby curve from $N_1$ to $N_2$ — impossible for a
  shortest geodesic unless $c$ is trivial, i.e.\ $N_1\cap N_2\ne\emptyset$.
  Identifying $N_1\cap N_2=E^{-1}(0)\subset\Omega_{N_1,N_2}(M)$, the same
  index bound shows $N_1\cap N_2\subset\Omega_{N_1,N_2}(M)$ is
  $(n-k_1-k_2)$-connected. Since $N_1\subset M$ is $(n-2k_1+1)$-connected by
  part (a) and $k_1\le k_2$ (so $n-2k_1+1\ge n-k_1-k_2+1>n-k_1-k_2$), the pair
  $\Omega_{N_1,N_2}(M)\subset\Omega_{M,N_2}(M)$ is also
  $(n-k_1-k_2)$-connected, hence so is
  $N_1\cap N_2\subset\Omega_{M,N_2}(M)$; since $\Omega_{M,N_2}(M)$ retracts
  onto $N_2$, this is $(n-k_1-k_2)$-connectivity of $N_1\cap N_2\to N_2$.
  \source{petersen:page-0281}
\end{proof}

\section{Further Study}

\begin{remark}[Frankel's theorem]
  \label{rem:pet-ch6-frankel-theorem}
  \lean{PetersenLib.frankelTheoremRemark}
  \uses{lem:pet-ch6-wilking-connectedness-principle}
  Part (b) of \cref{lem:pet-ch6-wilking-connectedness-principle}, in the
  special case that only nonemptiness of $N_1\cap N_2$ is retained, is
  classically known as \emph{Frankel's theorem}: two closed totally geodesic
  submanifolds $N_1,N_2$ of a manifold of positive curvature with
  $\dim N_1+\dim N_2\ge\dim M$ must intersect.
  \source{petersen:page-0282}
\end{remark}

\section{Exercises}

\begin{remark}[Exercise 6.7.1]
  \label{rem:pet-ch6-ex-6-7-1}
  \lean{PetersenLib.exercise_6_7_1}
  Show that in even dimensions the sphere and real projective space are the
  only closed manifolds with constant positive curvature.
  \source{petersen:page-0282}
\end{remark}

\begin{remark}[Exercise 6.7.2]
  \label{rem:pet-ch6-ex-6-7-2}
  \lean{PetersenLib.exercise_6_7_2}
  \uses{def:pet-ch6-parallel-field}
  Consider a rotationally symmetric metric $dr^2+\rho^2(r)d\theta^2$ and the
  cone $dr^2+(\rho(a)+\dot\rho(a)(r-a))^2d\theta^2$ tangent to it at the
  latitude $r=a$ (a genuine tangent cone when the surface is a surface of
  revolution). (1) Show that $\nabla_{\partial_\theta}$ agrees along $r=a$ on
  the two surfaces in the coordinates $(r,\theta)$, so parallel translation
  along $r=a$ is the same on both. (2) Unwrapping the cone to the flat plane,
  observe (by wrapping back up) that parallel translation along a latitude on
  a cone need not close up into a closed parallel field. (3) Show that the
  parallel field along $r=a$ closes up precisely when $\dot\rho(a)=0$.
  \source{petersen:page-0282}
\end{remark}

\begin{remark}[Exercise 6.7.3 — Fermi--Walker transport]
  \label{rem:pet-ch6-ex-6-7-3}
  \lean{PetersenLib.exercise_6_7_3}
  \uses{def:pet-ch6-vector-field-along-curve,def:pet-ch6-parallel-field}
  Let $c:[a,b]\to M$ have nonvanishing speed and $T=\dot c/|\dot c|$. Call $V$
  along $c$ a \emph{Fermi--Walker field} if
  $\dot V=g(V,T)\dot T-g(V,\dot T)T=(T\wedge\dot T)(V)$. (1) Show that given
  $V(t_0)$ there is a unique Fermi--Walker field with that value. (2) Show
  $T$ itself is a Fermi--Walker field. (3) Show $g(V,W)$ is constant along
  $c$ for Fermi--Walker fields $V,W$. (4) Show that if $c$ is a geodesic,
  Fermi--Walker fields are parallel.
  \source{petersen:page-0283}
\end{remark}

\begin{remark}[Exercise 6.7.4]
  \label{rem:pet-ch6-ex-6-7-4}
  \lean{PetersenLib.exercise_6_7_4}
  \uses{thm:pet-ch6-conjugate-radius-bound}
  Let $(M,g)$ be complete of constant curvature $k$ and $L:T_pM\to
  T_pS^n_k$ a linear isometry. If $k\le0$, show
  $\exp_p\circ L^{-1}\circ\exp_p^{-1}:S^n_k\to M$ is a Riemannian covering
  map. If $k>0$, show
  $\exp_p\circ L^{-1}\circ\exp_p^{-1}:S^n_k-\{-\tilde p\}\to M$ extends to a
  Riemannian covering map $S^n_k\to M$.
  \source{petersen:page-0283}
\end{remark}

\begin{remark}[Exercise 6.7.5]
  \label{rem:pet-ch6-ex-6-7-5}
  \lean{PetersenLib.exercise_6_7_5}
  \uses{def:pet-ch6-parallel-field}
  Let $c(s,t):[0,1]^2\to(M,g)$ satisfy
  $R(\partial c/\partial s,\partial c/\partial t)=0$. Show that for each
  $v\in T_{c(0,0)}M$ there is a parallel field $V:[0,1]^2\to TM$ along $c$
  (i.e.\ $\partial V/\partial s=\partial V/\partial t=0$ everywhere) with
  $V(0,0)=v$.
  \source{petersen:page-0283}
\end{remark}

\begin{remark}[Exercise 6.7.6]
  \label{rem:pet-ch6-ex-6-7-6}
  \lean{PetersenLib.exercise_6_7_6}
  \uses{lem:pet-ch6-third-partial-commutator}
  Using
  $R(\partial c/\partial s,\partial c/\partial t)\partial c/\partial u
  =\partial^3c/\partial s\,\partial t\,\partial u
  -\partial^3c/\partial t\,\partial s\,\partial u$, show that the two
  skew-symmetries and the first Bianchi identity for the curvature tensor
  (proposition 3.1.1 of the source) hold.
  \source{petersen:page-0283}
\end{remark}

\begin{remark}[Exercise 6.7.7]
  \label{rem:pet-ch6-ex-6-7-7}
  \lean{PetersenLib.exercise_6_7_7}
  \uses{def:pet-ch6-jacobi-field}
  Let $c$ be a geodesic and $X$ a Killing field. Show $X|_c$ is a Jacobi
  field.
  \source{petersen:page-0284}
\end{remark}

\begin{remark}[Exercise 6.7.8]
  \label{rem:pet-ch6-ex-6-7-8}
  \lean{PetersenLib.exercise_6_7_8}
  \uses{def:pet-ch6-jacobi-field,rem:pet-ch6-jacobi-dexp-relation}
  Let $c:[0,1]\to M$ be a geodesic. Show $\exp_{c(0)}$ has a critical point at
  $\dot c(0)$ iff there is a nontrivial Jacobi field $J$ along $c$ with
  $J(0)=0$, $\dot J(0)\perp\dot c(0)$, $J(1)=0$.
  \source{petersen:page-0284}
\end{remark}

\begin{remark}[Exercise 6.7.9]
  \label{rem:pet-ch6-ex-6-7-9}
  \lean{PetersenLib.exercise_6_7_9}
  \uses{def:pet-ch6-jacobi-field,lem:pet-ch6-hessian-lower-bound}
  Fix $p$, $v\in\operatorname{seg}_p^0$, $c(t)=\exp_p(tv)$, and the geodesic variation
  $\bar c(s,t)=\exp_p(t(v+sw))$ with Jacobi field $J$. Show
  $\nabla f_0|_{c(1)}=\dot c(1)$ and
  $\operatorname{Hess}f_0(J(1),J(1))=g(\dot J(1),\dot J(1))$ for $f_0=\tfrac12r^2$; use
  this to prove \cref{lem:pet-ch6-hessian-lower-bound} without first
  estimating $\operatorname{Hess}r$.
  \source{petersen:page-0284}
\end{remark}

\begin{remark}[Exercise 6.7.10]
  \label{rem:pet-ch6-ex-6-7-10}
  \lean{PetersenLib.exercise_6_7_10}
  \leanok
  \uses{def:pet-ch6-jacobi-field}
  Let $c$ be a geodesic and $J_1,J_2$ Jacobi fields along $c$. (1) Show
  $g(\dot J_1,J_2)-g(J_1,\dot J_2)$ is constant. (2) Show
  $g(J_1(t),\dot c(t))=g(J_1(0),\dot c(0))+g(\dot J_1(0),\dot c(0))t$.
  \source{petersen:page-0284}
\end{remark}

\begin{remark}[Exercise 6.7.11]
  \label{rem:pet-ch6-ex-6-7-11}
  \lean{PetersenLib.exercise_6_7_11}
  \uses{def:pet-ch6-jacobi-field}
  Let $J$ be a nontrivial Jacobi field along a unit speed geodesic $c$ with
  $J(0)=0$, $\dot J(0)\perp\dot c(0)$, and suppose $\sec\le K$. (1) With
  $\lambda=|J|^2/g(J,\dot J)$, show $\dot\lambda\le1+K\lambda^2$, $\lambda(0)=0$
  while $\lambda$ is defined. (2) Show that if $J(b)=0$ for some $b>0$, then
  $g(J(t),\dot J(t))=0$ for some $t\in(0,b)$; give an explicit example.
  \source{petersen:page-0284}
\end{remark}

\begin{remark}[Exercise 6.7.12]
  \label{rem:pet-ch6-ex-6-7-12}
  \lean{PetersenLib.exercise_6_7_12}
  \uses{def:pet-ch6-jacobi-field}
  Let $\mathfrak J$ be a self-adjoint space of Jacobi fields along a geodesic
  $c$ ($g(\dot J_1,J_2)=g(J_1,\dot J_2)$ for $J_1,J_2\in\mathfrak J$), and
  $\tilde{\mathfrak J}(t)=\{J(t):J\in\mathfrak J,\,J(t)=0\}
  +\{J(t):J\in\mathfrak J\}\subset T_{c(t)}M$ — read as the ambient span,
  with the two summands orthogonal. (1) Show the two summands are orthogonal.
  (2) Show $\{J\in\mathfrak J:J(t)=0\}$ is naturally isomorphic to the first
  summand. (3) Show $\dim\tilde{\mathfrak J}=\dim\tilde{\mathfrak J}(t)$ for
  all $t$.
  \source{petersen:page-0284}
\end{remark}

\begin{remark}[Exercise 6.7.13]
  \label{rem:pet-ch6-ex-6-7-13}
  \lean{PetersenLib.exercise_6_7_13}
  \uses{def:pet-ch6-vector-field-along-curve}
  Call $(M,g)$ \emph{$k$-point homogeneous} if for all $(p_1,\dots,p_k)$,
  $(q_1,\dots,q_k)$ with $|p_ip_j|=|q_iq_j|$ there is an isometry $F$ with
  $F(p_i)=q_i$ ($k=1$: homogeneous). (1) Show a homogeneous space has
  constant scalar curvature. (2) Show $k$-point homogeneous
  ($k>1$) implies $(k-1)$-point homogeneous. (3) Show two-point homogeneous
  implies Einstein. (4) Show three-point homogeneous implies constant
  curvature. (5) Show $\mathbb{RP}^2$ is not three-point homogeneous by
  exhibiting two noncongruent equilateral triangles of side length $\pi/3$.
  \source{petersen:page-0285}
\end{remark}

\begin{remark}[Exercise 6.7.14]
  \label{rem:pet-ch6-ex-6-7-14}
  \lean{PetersenLib.exercise_6_7_14}
  Starting with a geodesic on a two-dimensional space form, discuss how the
  equidistant curves change as they move away from the geodesic.
  \source{petersen:page-0285}
\end{remark}

\begin{remark}[Exercise 6.7.15]
  \label{rem:pet-ch6-ex-6-7-15}
  \lean{PetersenLib.exercise_6_7_15}
  \uses{def:pet-ch6-jacobi-field,thm:pet-ch6-rauch-comparison}
  Let $r(x)=|xp|$ with $-K\le\sec\le K$, $g=dr^2+g_r$ on
  $(0,R)\times S^{n-1}$, $2R<\operatorname{inj}_p$. (1) Show
  $\operatorname{sn}_K^2(r)\,ds_{n-1}^2\le g_r\le\operatorname{sn}_{-K}^2(r)\,ds_{n-1}^2$. (2) Show there
  is a universal constant $C$ with
  $|\operatorname{Hess}\tfrac12r^2-g|\le CK^2R^2$ for $R<\pi/(2\sqrt K)$.
  \source{petersen:page-0285}
\end{remark}

\begin{remark}[Exercise 6.7.16]
  \label{rem:pet-ch6-ex-6-7-16}
  \lean{PetersenLib.exercise_6_7_16}
  Let $(M,g)$ be complete. Show every element of $\pi_1(M,p)$ contains a
  shortest loop at $p$, and that this shortest loop is a geodesic loop.
  \source{petersen:page-0285}
\end{remark}

\begin{remark}[Exercise 6.7.17]
  \label{rem:pet-ch6-ex-6-7-17}
  \lean{PetersenLib.exercise_6_7_17}
  \uses{def:pet-ch6-path-space-AB}
  Let $(M,g)$ be complete with $\operatorname{inj}_p<R$ where
  $\exp_p:B(0,R)\to B(p,R)$ is nonsingular. Show the geodesic loop $c$ at $p$
  realizing the injectivity radius has index $0$.
  \source{petersen:page-0285}
\end{remark}

\begin{remark}[Exercise 6.7.18 — Frankel]
  \label{rem:pet-ch6-ex-6-7-18}
  \lean{PetersenLib.exercise_6_7_18}
  \uses{rem:pet-ch6-frankel-theorem}
  Let $M$ be $n$-dimensional with $\sec>0$ and $A,B$ closed totally geodesic
  submanifolds. Show directly (via a shortest connecting segment and the
  second variation formula, without invoking
  \cref{lem:pet-ch6-wilking-connectedness-principle}) that $A\cap B\ne
  \emptyset$ if $\dim A+\dim B\ge n$.
  \source{petersen:page-0286}
\end{remark}

\begin{remark}[Exercise 6.7.19]
  \label{rem:pet-ch6-ex-6-7-19}
  \lean{PetersenLib.exercise_6_7_19}
  \uses{def:pet-ch6-path-space-AB,lem:pet-ch6-wilking-connectedness-principle}
  Let $M$ be complete, $A\subset M$ compact, without invoking
  \cref{lem:pet-ch6-wilking-connectedness-principle}. (1) Show non-stationary
  curves in $\Omega_{A,A}(M)$ deform to shorter curves in $\Omega_{A,A}(M)$.
  (2) Show stationary curves are geodesics perpendicular to $A$ at the
  endpoints. (3) If $M$ has $\sec>0$, $A$ is totally geodesic, and $2\dim
  A\ge\dim M$, show all stationary curves deform to shorter curves in
  $\Omega_{A,A}(M)$. (4) (Wilking) Conclude that $\pi_1(M,A)$ is trivial.
  \source{petersen:page-0286}
\end{remark}

\begin{remark}[Exercise 6.7.20]
  \label{rem:pet-ch6-ex-6-7-20}
  \lean{PetersenLib.exercise_6_7_20}
  \uses{thm:pet-ch6-preissmann,def:pet-ch6-axis-displacement}
  Generalize \cref{thm:pet-ch6-preissmann} to show any solvable subgroup of
  $\pi_1$ of a compact negatively curved manifold is cyclic. Hint: use
  torsion-freeness and, by contradiction and solvability, a subgroup
  generated by deck transformations $F,G$ with $F\circ G=G^k\circ F$,
  $k\ne0$; show that if $c$ is an axis for $G$, $F\circ c$ is an axis for
  $G^k$, and use uniqueness of axes for $G^k$ for a contradiction.
  \source{petersen:page-0286}
\end{remark}

\begin{remark}[Exercise 6.7.21]
  \label{rem:pet-ch6-ex-6-7-21}
  \lean{PetersenLib.exercise_6_7_21}
  \uses{rem:pet-ch6-closed-parallel-field-construction}
  Let $(M,g)$ be compact with $\sec>0$ and $F:M\to M$ an isometry of finite
  order without fixed points. Show that if $\dim M$ is even, $F$ is
  orientation reversing; if odd, orientation preserving. (Weinstein: this
  holds even without the finite-order hypothesis.)
  \source{petersen:page-0286}
\end{remark}

\begin{remark}[Exercise 6.7.22]
  \label{rem:pet-ch6-ex-6-7-22}
  \lean{PetersenLib.exercise_6_7_22}
  \uses{thm:pet-ch6-cartan-fixed-point}
  Use an analogue of \cref{thm:pet-ch6-cartan-fixed-point} to show a closed
  manifold of constant curvature $1$ is either the standard sphere or has
  diameter $\le\pi/2$; generalize to show a closed manifold with $\sec\ge1$
  is either simply connected or has diameter $\le\pi/2$.
  \source{petersen:page-0286}
\end{remark}

\begin{remark}[Exercise 6.7.23]
  \label{rem:pet-ch6-ex-6-7-23}
  \lean{PetersenLib.exercise_6_7_23}
  \uses{thm:pet-ch6-rauch-comparison}
  Let $(M,g)$ be complete $n$-dimensional with $|\sec|<K$, and fix $p$,
  $\epsilon>0$, points $p_i$ with $r^i(x)=|xp_i|$ smooth on $B(p,\epsilon)$,
  $g(\nabla r^i,\nabla r^j)|_p=\delta^{ij}$, $|pp_i|\ge2\epsilon$. Let
  $g^{ij}=g(\nabla r^i,\nabla r^j)$. (1) Show there is $C=C(n,K,\epsilon)$
  with $|dg^{ij}|\le C$. (2) Show $C$ can be chosen so that
  $C(n,\lambda^{-2}K,\lambda\epsilon)\to0$ as $\lambda\to\infty$. (3) Show
  there is $\delta=\delta(n,C)>0$ with $g^{ij}$ invertible on $B(p,\delta)$
  and its inverse $g_{ij}$ satisfying $|g_{ij}-\delta_{ij}|\le\tfrac12$,
  $|dg_{ij}|\le C'(n,K,\epsilon)$. (4) Show
  $(r^1(x)-r^1(p),\dots,r^n(x)-r^n(p))$ is a coordinate system on
  $B(p,\delta)$ whose image contains $B(0,\delta/2)$.
  \source{petersen:page-0286,petersen:page-0287}
\end{remark}

\begin{remark}[Exercise 6.7.24 — the index form]
  \label{rem:pet-ch6-ex-6-7-24}
  \lean{PetersenLib.exercise_6_7_24,PetersenLib.indexForm}
  \uses{rem:pet-ch6-second-variation-special-cases,def:pet-ch6-jacobi-field}
  For $V,W$ along a geodesic $c:[0,1]\to(M,g)$, the \emph{index form} is
  \[
    I_0^1(V,W)=I(V,W)=\int_0^1\big(g(\dot V,\dot W)-g(R(V,\dot c)\dot c,W)\big)\,dt,
  \]
  agreeing with the second variation of energy when $V,W$ come from a proper
  variation. Assume $c$ locally minimizes energy, so $I(V,V)\ge0$ for proper
  variational fields. (1) If $I(V,V)=0$ for such $V$, show $V$ is a Jacobi
  field (expand $I(V+\epsilon W,V+\epsilon W)\ge0$ in $\epsilon$ for all
  proper $W$). (2) If $V(0)=J(0)$, $V(1)=J(1)$ and $J$ is a Jacobi field,
  show $I(V,J)=I(J,J)$. (3) \emph{(Index lemma)} If in addition no nontrivial
  Jacobi field vanishes at both endpoints of $c$, show $I(V,V)\ge I(J,J)$
  with equality iff $V=J$ (via $0\le I(V-J,V-J)=I(V,V)-I(J,J)$). (4) If a
  nontrivial Jacobi field $J$ vanishes at $0$ and $1$, show
  $c:[0,1+\epsilon]\to M$ is not locally minimizing for $\epsilon>0$, by
  patching $J$ on $[0,1-\epsilon]$ to a Jacobi field $K$ on $[1-\epsilon,
  1+\epsilon]$ with $K(1-\epsilon)=J(1-\epsilon)$, $K(1+\epsilon)=0$, and
  comparing $I_0^1(J,J)=0$ against $I_0^{1-\epsilon}(J,J)+I_{1-\epsilon}^{1+
  \epsilon}(K,K)$, which is strictly smaller since $K\ne J$ on
  $[1-\epsilon,1]$.
  \source{petersen:page-0287,petersen:page-0288}
\end{remark}

\begin{remark}[Exercise 6.7.25 — index comparison]
  \label{rem:pet-ch6-ex-6-7-25}
  \lean{PetersenLib.exercise_6_7_25}
  \uses{rem:pet-ch6-ex-6-7-24,thm:pet-ch6-rauch-comparison}
  Let $J$ be a nontrivial Jacobi field along a unit speed geodesic $c$ with
  $J(0)=0$, $\dot J(0)\perp\dot c(0)$, on a manifold with $\sec\ge k$, and
  suppose no Jacobi field on $c|_{[0,b]}$ vanishes at both endpoints. (1)
  Show $I_0^b(J,J)=g(J(b),\dot J(b))$. (2) With $V(t)=\frac{\operatorname{sn}_k(t)}
  {\operatorname{sn}_k(b)}E(t)$ for a parallel field $E$ with $E(b)=J(b)$, show
  $I_0^b(V,V)\le\frac{\operatorname{sn}_k'(b)}{\operatorname{sn}_k(b)}|J(b)|^2$. (3) Conclude
  $g(J(b),\dot J(b))\le\frac{\operatorname{sn}_k'(b)}{\operatorname{sn}_k(b)}|J(b)|^2$, and use this to
  prove the lower-curvature-bound half of
  \cref{thm:pet-ch6-rauch-comparison}.
  \source{petersen:page-0288}
\end{remark}

\begin{remark}[Exercise 6.7.26]
  \label{rem:pet-ch6-ex-6-7-26}
  \lean{PetersenLib.exercise_6_7_26}
  \uses{thm:pet-ch6-cartan-fixed-point,def:pet-ch6-convexity-radius}
  Let $\mathbf G\subset\operatorname{Iso}(M,g)$, with $\operatorname{Iso}(M,g)$ given the
  compact-open topology. (1) If $M$ is complete, simply connected,
  nonpositively curved, and $\mathbf G$ is compact, show $\mathbf G$ has a
  fixed point (imitate \cref{thm:pet-ch6-cartan-fixed-point}). (2) Call
  $\mathbf G$ \emph{$(p,\epsilon)$-small} if $\mathbf Gp\subset\bar
  B(p,\epsilon)$; show that for sufficiently small $\epsilon(p)$ the closure
  $\bar{\mathbf G}$ of a $(p,\epsilon)$-small group is compact and
  $(p,\epsilon)$-small (without assuming $M$ complete). (3) Show a
  $(p,\epsilon)$-small $\mathbf G$ is $(q,2|pq|+\epsilon)$-small. (4) If
  $\epsilon$ is much smaller than the convexity radius on $\bar
  B(p,4\epsilon)$, show $\mathbf G$ has a fixed point. (5) Show that for
  every $p$ there is $\epsilon>0$ such that no subgroup of $\operatorname{Iso}(M,g)$ can
  be $(p,\epsilon)$-small (by making $\mathbf G$ act freely on a suitable
  subset of a product $M\times\cdots\times M$). (6) Conclude
  $\operatorname{Iso}(M,g)$ has no small subgroups (no neighborhood of the identity
  contains a nontrivial subgroup).
  \source{petersen:page-0288,petersen:page-0289}
\end{remark}

\begin{remark}[Exercise 6.7.27]
  \label{rem:pet-ch6-ex-6-7-27}
  \lean{PetersenLib.exercise_6_7_27}
  \uses{def:pet-ch6-parallel-field}
  Construct a Riemannian metric on $TM$ for a Riemannian manifold $(M,g)$
  such that $\pi:TM\to M$ is a Riemannian submersion restricting to the given
  Euclidean metric on each fiber, by declaring parallel fields along curves
  in $M$ to give horizontal lifts.
  \source{petersen:page-0289}
\end{remark}

\begin{remark}[Exercise 6.7.28]
  \label{rem:pet-ch6-ex-6-7-28}
  \lean{PetersenLib.exercise_6_7_28}
  \uses{def:pet-ch6-parallel-field}
  For a Riemannian manifold $(M,g)$ let $FM\to M$ be the orthonormal frame
  bundle. Find a Riemannian metric on $FM$ making $\pi:FM\to M$ a Riemannian
  submersion with fibers isometric to $O(n)$, by declaring parallel
  orthonormal frames along curves in $M$ to give horizontal lifts.
  \source{petersen:page-0289}
\end{remark}
